%Version 3.1 December 2024
% See section 11 of the User Manual for version history
%
%%%%%%%%%%%%%%%%%%%%%%%%%%%%%%%%%%%%%%%%%%%%%%%%%%%%%%%%%%%%%%%%%%%%%%
%%                                                                 %%
%% Please do not use \input{...} to include other tex files.       %%
%% Submit your LaTeX manuscript as one .tex document.              %%
%%                                                                 %%
%% All additional figures and files should be attached             %%
%% separately and not embedded in the \TeX\ document itself.       %%
%%                                                                 %%
%%%%%%%%%%%%%%%%%%%%%%%%%%%%%%%%%%%%%%%%%%%%%%%%%%%%%%%%%%%%%%%%%%%%%

%%\documentclass[referee,sn-basic]{sn-jnl}% referee option is meant for double line spacing

%%=======================================================%%
%% to print line numbers in the margin use lineno option %%
%%=======================================================%%

%%\documentclass[lineno,pdflatex,sn-basic]{sn-jnl}% Basic Springer Nature Reference Style/Chemistry Reference Style

%%=========================================================================================%%
%% the documentclass is set to pdflatex as default. You can delete it if not appropriate.  %%
%%=========================================================================================%%

%%\documentclass[sn-basic]{sn-jnl}% Basic Springer Nature Reference Style/Chemistry Reference Style

%%Note: the following reference styles support Namedate and Numbered referencing. By default the style follows the most common style. To switch between the options you can add or remove “Numbered” in the optional parenthesis. 
%%The option is available for: sn-basic.bst, sn-chicago.bst%  
 
%%\documentclass[pdflatex,sn-nature]{sn-jnl}% Style for submissions to Nature Portfolio journals
%%\documentclass[pdflatex,sn-basic]{sn-jnl}% Basic Springer Nature Reference Style/Chemistry Reference Style
\documentclass[pdflatex,sn-mathphys-num]{sn-jnl}% Math and Physical Sciences Numbered Reference Style
%%\documentclass[pdflatex,sn-mathphys-ay]{sn-jnl}% Math and Physical Sciences Author Year Reference Style
%%\documentclass[pdflatex,sn-aps]{sn-jnl}% American Physical Society (APS) Reference Style
%%\documentclass[pdflatex,sn-vancouver-num]{sn-jnl}% Vancouver Numbered Reference Style
%%\documentclass[pdflatex,sn-vancouver-ay]{sn-jnl}% Vancouver Author Year Reference Style
%%\documentclass[pdflatex,sn-apa]{sn-jnl}% APA Reference Style
%%\documentclass[pdflatex,sn-chicago]{sn-jnl}% Chicago-based Humanities Reference Style

%%%% Standard Packages
%%<additional latex packages if required can be included here>

\usepackage{graphicx}%
\usepackage{multirow}%
\usepackage{amsmath,amssymb,amsfonts}%
\usepackage{amsthm}%
\usepackage{mathrsfs}%
\usepackage[title]{appendix}%
\usepackage{xcolor}%
\usepackage{textcomp}%
\usepackage{manyfoot}%
\usepackage{booktabs}%
\usepackage{algorithm}%
\usepackage{algorithmicx}%
\usepackage{algpseudocode}%
\usepackage{listings}%
%%%%

%%%%%=============================================================================%%%%
%%%%  Remarks: This template is provided to aid authors with the preparation
%%%%  of original research articles intended for submission to journals published 
%%%%  by Springer Nature. The guidance has been prepared in partnership with 
%%%%  production teams to conform to Springer Nature technical requirements. 
%%%%  Editorial and presentation requirements differ among journal portfolios and 
%%%%  research disciplines. You may find sections in this template are irrelevant 
%%%%  to your work and are empowered to omit any such section if allowed by the 
%%%%  journal you intend to submit to. The submission guidelines and policies 
%%%%  of the journal take precedence. A detailed User Manual is available in the 
%%%%  template package for technical guidance.
%%%%%=============================================================================%%%%

%% as per the requirement new theorem styles can be included as shown below
\theoremstyle{thmstyleone}%
\newtheorem{theorem}{Theorem}%  meant for continuous numbers
%%\newtheorem{theorem}{Theorem}[section]% meant for sectionwise numbers
%% optional argument [theorem] produces theorem numbering sequence instead of independent numbers for Proposition
\newtheorem{proposition}[theorem]{Proposition}% 
%%\newtheorem{proposition}{Proposition}% to get separate numbers for theorem and proposition etc.

\theoremstyle{thmstyletwo}%
\newtheorem{example}{Example}%
\newtheorem{remark}{Remark}%

\theoremstyle{thmstylethree}%
\newtheorem{definition}{Definition}%

\raggedbottom
%%\unnumbered% uncomment this for unnumbered level heads

\begin{document}

\title[Article Title]{Reinforcement Learning for Within-Visual-Range Air Combat with Curriculum Learning and Dynamic Opponent Pooling}

%%=============================================================%%
%% GivenName	-> \fnm{Joergen W.}
%% Particle	-> \spfx{van der} -> surname prefix
%% FamilyName	-> \sur{Ploeg}
%% Suffix	-> \sfx{IV}
%% \author*[1,2]{\fnm{Joergen W.} \spfx{van der} \sur{Ploeg} 
%%  \sfx{IV}}\email{iauthor@gmail.com}
%%=============================================================%%

% \author*[1,2]{\fnm{First} \sur{Author}}\email{iauthor@gmail.com}

% \author[2,3]{\fnm{Second} \sur{Author}}\email{iiauthor@gmail.com}
% \equalcont{These authors contributed equally to this work.}

% \author[1,2]{\fnm{Third} \sur{Author}}\email{iiiauthor@gmail.com}
% \equalcont{These authors contributed equally to this work.}

% \affil*[1]{\orgdiv{Department}, \orgname{Organization}, \orgaddress{\street{Street}, \city{City}, \postcode{100190}, \state{State}, \country{Country}}}

% \affil[2]{\orgdiv{Department}, \orgname{Organization}, \orgaddress{\street{Street}, \city{City}, \postcode{10587}, \state{State}, \country{Country}}}

% \affil[3]{\orgdiv{Department}, \orgname{Organization}, \orgaddress{\street{Street}, \city{City}, \postcode{610101}, \state{State}, \country{Country}}}

% Author information
\author[1]{\fnm{Yue} \sur{Yu}}

\author[1]{\fnm{Yuhang} \sur{Wang}}

\author[2]{\fnm{Zhe} \sur{Zhang}}

\author[1]{\fnm{Bei} \sur{Lu}}

\author*[1]{\fnm{Qifu} \sur{Li}}\email{qifuli@sjtu.edu.cn}

\affil*[1]{\orgdiv{School of Aeronautics and Astronautics},
\orgname{Shanghai Jiao Tong University},
\orgaddress{\city{Shanghai}, \country{China}}}

\affil[2]{\orgname{Chinese Flight Test Establishment}, \orgaddress{\city{Xi’an}, \country{China}}}

%%==================================%%
%% Sample for unstructured abstract %%
%%==================================%%

% \abstract{Achieving stable training and robust generalization in autonomous air combat remains a significant challenge for unmanned combat aircraft. This study proposes a reinforcement learning framework for within-visual-range dogfights. Our approach first establishes a six-degree-of-freedom air combat simulation environment. Then, a proximal policy optimization based architecture, featuring a curriculum learning and a dynamic opponent pool, is designed to mitigate sparse rewards and policy overfitting during self-play. Engagement simulations indicate that the proposed framework improves tactical performance, demonstrating stable convergence and competitive win rates.}
\abstract{
Achieving stable training and robust generalization in autonomous air combat remains a significant challenge for unmanned combat aircraft. This study proposes a reinforcement learning framework for within-visual-range dogfights in a six-degree-of-freedom air combat simulation environment based on JSBSim. A proximal policy optimization-based architecture is developed, featuring curriculum learning and a dynamic opponent pool. The curriculum progressively shifts the learning objective from stable flight control to tactical air-combat behavior, while gradually increasing opponent threat levels to alleviate sparse-reward exploration difficulty. The dynamic opponent pool combines protected opponents with competitive historical policy snapshots, improving adversarial diversity and mitigating policy overfitting during self-play. Engagement simulations demonstrate that the proposed framework achieves stable convergence and improved tactical performance, with the best proposed model improving the win rate over the baseline by 55.2 percentage points.
}

%%================================%%
%% Sample for structured abstract %%
%%================================%%

% \abstract{\textbf{Purpose:} The abstract serves both as a general introduction to the topic and as a brief, non-technical summary of the main results and their implications. The abstract must not include subheadings (unless expressly permitted in the journal's Instructions to Authors), equations or citations. As a guide the abstract should not exceed 200 words. Most journals do not set a hard limit however authors are advised to check the author instructions for the journal they are submitting to.
% 
% \textbf{Methods:} The abstract serves both as a general introduction to the topic and as a brief, non-technical summary of the main results and their implications. The abstract must not include subheadings (unless expressly permitted in the journal's Instructions to Authors), equations or citations. As a guide the abstract should not exceed 200 words. Most journals do not set a hard limit however authors are advised to check the author instructions for the journal they are submitting to.
% 
% \textbf{Results:} The abstract serves both as a general introduction to the topic and as a brief, non-technical summary of the main results and their implications. The abstract must not include subheadings (unless expressly permitted in the journal's Instructions to Authors), equations or citations. As a guide the abstract should not exceed 200 words. Most journals do not set a hard limit however authors are advised to check the author instructions for the journal they are submitting to.
% 
% \textbf{Conclusion:} The abstract serves both as a general introduction to the topic and as a brief, non-technical summary of the main results and their implications. The abstract must not include subheadings (unless expressly permitted in the journal's Instructions to Authors), equations or citations. As a guide the abstract should not exceed 200 words. Most journals do not set a hard limit however authors are advised to check the author instructions for the journal they are submitting to.}

\keywords{reinforcement learning, within-visual-range air combat, dogfight, curriculum learning, dynamic opponent pool}

%%\pacs[JEL Classification]{D8, H51}

%%\pacs[MSC Classification]{35A01, 65L10, 65L12, 65L20, 65L70}

\maketitle

\section{Introduction}\label{sec1}

Recent advances in artificial intelligence have driven air combat away from rule-based systems and toward data-driven approaches \cite{Strand_2025}. Dogfight scenarios are especially challenging because they require aircraft to interpret high-dimensional states and generate continuous control commands at high frequency under rapidly changing conditions. Reinforcement learning (RL) improves behavior through repeated trial-and-error in the environment, making it well suited to dynamic adversarial settings. The potential of RL for dogfight has been supported by recent studies. A representative example is the 2020 Defense Advanced Research Projects Agency (DARPA) AlphaDogfight Trials (ADT), in which an artificial intelligence (AI) agent defeated an experienced human pilot in a one-on-one within-visual-range (WVR) dogfight \cite{pope_hierarchical_2023}. This result suggested that AI could achieve competitive performance in tactical decision-making and maneuvering control, and it further stimulated research on autonomous unmanned combat aircraft.


Even so, robust performance in dogfight scenarios remains difficult to achieve. Specifically, the three major challenges are the training instability inherent in high-dimensional continuous control, the exploration difficulty under sparse rewards, and the policy overfitting caused by self-play.

The high-dimensional and continuous nature of aircraft dynamics often leads to convergence difficulty during training. To alleviate this, many studies have adopted hierarchical architectures to decouple complex control tasks. For instance, at the low-level maneuvering layer, RL-based controllers have outperformed conventional proportional-integral-derivative (PID) controllers in response speed and tracking accuracy, while also showing stronger adaptability to severe wind disturbances and aggressive maneuvers \cite{koch_reinforcement_2019, bohn_data-efficient_2023, oconnell_neural-fly_2022, clarke_deep_2020, clarke_reinforcement_2023}. Similarly, to address the curse of dimensionality and the credit assignment problem in multi-aircraft cooperation, hierarchical architectures and value decomposition methods have been used \cite{chai_hierarchical_2022, wang_hierarchical_2023, feng_autonomous_2023, selmonaj_hierarchical_2023}. Methods including hierarchical control and gain scheduling have also been proposed to address convergence difficulty and oscillatory behavior in end-to-end control \cite{chen_physics-informed_2026}. Notably, one study based on JSBSim, an open-source flight dynamics model widely used for aircraft simulation \cite{jsbsim}, demonstrated the effectiveness of a hierarchical RL architecture in experiments against a human pilot \cite{pope_hierarchical_2023}. While hierarchical methods simplify the learning process, they often require manual design and may introduce sub-optimalities at the interfaces between decision-making and control layers. Consequently, end-to-end control frameworks are desirable for their potential to learn globally optimal behaviors, though achieving stable convergence remains challenging. To improve end-to-end training stability and sample efficiency, researchers have proposed progressive multi-objective strategy optimization methods \cite{li_progressive_2025}, as well as model-based reinforcement learning \cite{GIRAL2026110985}. Furthermore, sequence models such as Transformers and long short-term memory networks (LSTMs) have been integrated to enhance long-horizon decision-making and state estimation under partial observability \cite{yu_hybrid_2025, yang_close-range_2025, hsu_scalable_2021, bae_deep_2023}. 

Beyond control stability, dogfight tasks typically suffer from exploration difficulties due to sparse reward signals. To alleviate this issue and improve RL performance, curriculum learning and imitation learning have been actively explored \cite{WANG2024109402, li_imitative_2025}. The performance of imitation learning remains limited. For instance, it has been observed that even when facing non-adversarial opponents performing basic maneuvers such as level flight, circling, or serpentine patterns, the shooting success rates achieved through imitation learning still fall short of those seen in expert demonstrations \cite{li_imitative_2025}. Studies employing curriculum learning progressively expose agents to increasingly challenging scenarios, such as gradually reducing initial tactical advantages or introducing more competitive opponents \cite{bae_deep_2023, kong_hierarchical, han_deep_2022, Selmonaj_2023}. Therefore, curriculum learning serves as an effective approach to smoothly guiding the exploration process and overcoming the sparse reward challenge in complex dogfight scenarios.

Within a reinforcement learning framework, improving generalization requires addressing policy overfitting and cyclic dominance that often arise during self-play \cite{han_augmenting_2025, chen_autonomous_2025}. To promote gradual and stable performance improvement, one study employed simulated annealing to design a dynamic opponent-pool update rule: a new policy was added to the opponent pool only when its win rate exceeded 50\%, the oldest policy was removed, and the sampling distribution was progressively shifted toward stronger opponents \cite{chen_autonomous_2025}. Another study approached the problem from the perspective of population diversity, introducing a population-entropy regularization term to enlarge the KL divergence among policies in the opponent pool \cite{han_augmenting_2025}. While this entropy-based method improved the generalization of the main agent, its validation was limited to a low-fidelity three-degree-of-freedom point-mass model with a discrete action space. Building on these insights, designing a robust dynamic opponent pool that effectively balances diversity and competence in high-fidelity, continuous-action environments remains essential. Furthermore, a challenge in such frameworks is the cold-start problem, which concerns the selection of appropriate initial opponents to kickstart the learning process. Whether utilizing expert strategies \cite{li_imitative_2025} or pre-trained models \cite{han_augmenting_2025} as baseline adversaries, the competence and diversity of these initial opponents can significantly influence the convergence stability and the ultimate proficiency of the agent.

While RL shows great promise for air combat, convergence under sparse rewards, cyclic dominance, and overfitting in self-play continue to be important open problems. These intertwined challenges form the central focus of the present study.

This work presents a reinforcement learning framework for one-on-one within-visual-range gun-based dogfights, and evaluates it in a six-degree-of-freedom (6-DOF) air combat environment built on JSBSim \cite{jsbsim}. The environment models an F-16 aircraft, and adopts engagement criteria consistent with the DARPA setting reported in \cite{pope_hierarchical_2023}. Proximal policy optimization (PPO) is used as the underlying learning algorithm. To address the sparse rewards and overfitting tendencies commonly encountered in self-play air combat training, the framework incorporates a curriculum learning framework and a dynamic opponent pool.

The main contributions of this work are summarized as follows:
\begin{itemize}
    \item A curriculum learning framework is developed, consisting of an objective curriculum that transitions from stable flight control to air-combat objectives, and an opponent curriculum in which adversary difficulty increases progressively over training.
    \item A dynamic opponent pool is introduced to maintain policy diversity through evaluation during training. In addition, a goalpoint-based RL opponent is incorporated to mitigate the cold-start problem.
\end{itemize}

The remainder of this paper is organized as follows. Section \ref{sec:env_setup} presents the mathematical formulation of the air combat task and the construction of the simulation environment. Section \ref{sec:RL_training} describes the PPO-based training framework, including the curriculum and the dynamic opponent-pool update strategy. Section \ref{sec:experiment} evaluates the tactical performance of the trained agent across multiple engagement scenarios. Section \ref{sec:conclusion} concludes the paper.

\section{Environment Setup}
\label{sec:env_setup}

\subsection{Air Combat Formulation}
\label{sec:mdp_modeling}

WVR air combat is a dynamic and nonlinear adversarial process. In the one-on-one WVR gun dogfight scenario considered in this study, two aircraft execute complex tactical maneuvers in a continuous three-dimensional space in order to secure a favorable firing position behind the opponent. Under a gun-only weapon configuration, the aircraft must avoid incoming fire while closing on the opponent and satisfying range constraints to achieve a valid kill. This setting requires the policy to handle long-horizon spatiotemporal interactions under rapidly changing engagement conditions.

Under a fixed opponent policy, the dogfight process can be approximated from the perspective of the agent as a Markov decision process (MDP), represented by the tuple $(\mathcal{S}, \mathcal{A}, \mathcal{P}, \mathcal{R})$. Here, the state space $\mathcal{S}$ and action space $\mathcal{A}$ define the interaction interface between the aircraft and the environment, while $\mathcal{P}$ denotes the transition function and $\mathcal{R}$ denotes the reward function.

The state space $\mathcal{S}$ provides the agent with a structured description of the current engagement and the aircraft's own flight state. In this work, the observation is represented as a continuous vector. Table \ref{tab:key_state_variables} summarizes the observation variables, while their full definitions and bounds are provided in Table \ref{tab:full_state_space}. The angles of relative geometry are shown in Figure~\ref{fig:track}.


% \begin{table}
%     \centering
%     \caption{Definitions of key state variables in observation space.}
%     \label{tab:key_state_variables}
%     \setlength{\tabcolsep}{4pt}
%     \renewcommand{\arraystretch}{1.05}
%     \begin{tabular}{>{\raggedright\arraybackslash}m{0.24\columnwidth}
%                     >{\raggedright\arraybackslash}m{0.66\columnwidth}}
%         \toprule
%         \textbf{Category} & \textbf{Variable Symbol \& Meaning} \\
%         \midrule
%         \multirow{3}{=}{Ego Basic States}
%             & $h$, Altitude \\
%             & $u, v, w$, Body-axis velocities \\
%             & $p, q, r$, Body-axis angular velocities \\
%             & $\theta, \phi$, Pitch and roll angles \\
%         \midrule
%         \multirow{4}{=}{Relative Geometry}
%             & $D$, Relative distance \\
%             & $V_c$, Closure rate \\
%             & $\theta_{tra}$, Track angle \\
%             & $\theta_{adv}$, Adverse angle \\
%             & $\theta_{bea}$, Bearing angle \\
%             & $\theta_{ele}$, Elevation angle \\
%         \midrule
%         Opponent Features
%             & $\mathbf{x}_{oppo}$, Ego basic states and relative geometry from the opponent's perspective \\
%         \midrule
%         \multirow{2}{=}{Action Variables}
%             & $\delta_a, \delta_e, \delta_r$, Aileron, elevator, and rudder commands \\
%             & $\delta_T$, Throttle command \\
%         \botrule
%     \end{tabular}
% \end{table}
\begin{table}[h]
\caption{Definitions of key state variables in the observation space.}\label{tab:key_state_variables}%
\begin{tabular*}{\textwidth}{@{\extracolsep\fill}p{0.28\textwidth}p{0.68\textwidth}}
\toprule
Category & Variable symbol and meaning \\
\midrule
\multirow{4}{=}{Ego Basic States}
    & $h$, Altitude \\
    & $u, v, w$, Body-axis velocities \\
    & $p, q, r$, Body-axis angular velocities \\
    & $\theta, \phi$, Pitch and roll angles \\
\midrule
\multirow{6}{=}{Relative Geometry}
    & $D$, Relative distance \\
    & $V_c$, Closure rate \\
    & $\theta_{tra}$, Track angle \\
    & $\theta_{adv}$, Adverse angle \\
    & $\theta_{bea}$, Bearing angle \\
    & $\theta_{ele}$, Elevation angle \\
\midrule
\multirow{2}{=}{Opponent Features}
    & $\mathbf{x}_{oppo}$, Ego basic states and relative geometry from the opponent's perspective \\
\midrule
\multirow{2}{=}{Action Variables}
    & $\delta_a, \delta_e, \delta_r$, Aileron, elevator, and rudder commands \\
    & $\delta_T$, Throttle command \\
\botrule
\end{tabular*}
\end{table}

% \begin{figure}
%     \centering
%     \includegraphics[width=0.49\linewidth]{graph/track_angle_finish.png}
%     \hfill
%     \includegraphics[width=0.49\linewidth]{graph/elevation_angle_finish.png}
%     \caption{Depiction of the track angle, adverse angle, elevation angle, and bearing angle.}
%     \label{fig:track}
% \end{figure}
\begin{figure}[h]
\centering
\includegraphics[width=0.9\textwidth]{graph/relative_geometry_angles.png}
\caption{Definition of angular variables in air combat geometry.}\label{fig:track}
\end{figure}

Instead of using a hierarchical control architecture based on high-level waypoints or attitude-command decomposition, this work adopts a mapping from observations directly to pilot inputs. Accordingly, the action space $\mathcal{A}$ is defined as a four-dimensional continuous vector, whose components correspond to normalized target roll rate, target pitch rate, target yaw rate, and throttle commands.

The transition function $\mathcal{P}(s_{t+1}\mid s_t, a_t)$ is governed by the equations of motion evaluated by JSBSim. Episode termination is jointly determined by weapon effectiveness, flight-safety constraints, and a maximum engagement duration. Specifically, an episode terminates when an aircraft falls within the opponent’s gun engagement zone and its health points (HP) are reduced to zero. In addition, if an aircraft crashes into the ground during aggressive maneuvering, that is, when the altitude satisfies $h \le 0$, the episode is forcibly terminated. Finally, to prevent prolonged circling at long range or static stalemates, the simulator imposes a maximum episode length of $60$ seconds of physical time. Once this limit is reached, the engagement is truncated and scored as a draw.


\subsection{Simulation Environment}
\label{sec:simulation_architecture}

This work constructs the training environment using JSBSim, an open-source high-fidelity flight dynamics engine. The open-source project Gym-jsbsim \cite{jsbgym} provides a viable approach to wrapping JSBSim into a Gymnasium-compatible environment. During simulation, the environment state is updated at a frequency of $60\text{ Hz}$.

The core dynamics of the system are governed by the real-time integration of a set of ordinary differential equations. The translational equations of motion can be written as
\begin{equation}
    \begin{cases}
        \dot{u} = r v - q w + \frac{1}{m}\left(W_x + F_x^{(A)} + F_x^{(T)}\right) \\
        \dot{v} = -r u + p w + \frac{1}{m}\left(W_y + F_y^{(A)} + F_y^{(T)}\right) \\
        \dot{w} = q u - p v + \frac{1}{m}\left(W_z + F_z^{(A)} + F_z^{(T)}\right)
    \end{cases},
    \label{eq:clmes}
\end{equation}
where $u, v, w$ and $p, q, r$ denote the translational velocities and angular rates along the three body axes, respectively; $m$ is the aircraft mass; $W_{x,y,z}$ are the gravity components; and $F^{(A)}$ and $F^{(T)}$ represent the resultant aerodynamic forces and engine thrust, respectively. Correspondingly, the rotational equations of motion are given by
\begin{equation}
    \begin{cases}
        \dot{p} = (C_1 r + C_2 p)q + C_3 \mathcal{L} + C_4 \mathcal{N} \\
        \dot{q} = C_5 p r - C_6 (p^2 - r^2) + C_7 \mathcal{M} \\
        \dot{r} = (C_8 p - C_2 r)q + C_4 \mathcal{L} + C_9 \mathcal{N}
    \end{cases},
    \label{eq:cames}
\end{equation}
where $\mathcal{L}$, $\mathcal{M}$, and $\mathcal{N}$ denote the total rolling, pitching, and yawing moments generated by aerodynamic and propulsive effects. The coefficients $C_1$ to $C_9$ are determined entirely by the rigid-body inertia properties of the aircraft, which are defined as follows.

Let $Ox_By_Bz_B$ denote the body-fixed frame, whose origin $O$ is located at the aircraft centre of gravity, with the $x_B$ axis pointing forward along the fuselage reference line, the $y_B$ axis pointing to the right wing, and the $z_B$ axis pointing downward. The inertia tensor of the aircraft expressed in this frame is
\begin{equation}
    \mathbf{J} =
    \begin{bmatrix}
         I_{xx} & -I_{xy} & -I_{xz} \\[2pt]
        -I_{xy} &  I_{yy} & -I_{yz} \\[2pt]
        -I_{xz} & -I_{yz} &  I_{zz}
    \end{bmatrix},
    \label{eq:inertia_tensor}
\end{equation}
where $(x,y,z)$ denote the coordinates of a mass element $\mathrm{d}m$ in the body frame, the diagonal entries
\begin{equation}
    I_{xx} = \int \left(y^{2} + z^{2}\right) \mathrm{d}m, \qquad
    I_{yy} = \int \left(x^{2} + z^{2}\right) \mathrm{d}m, \qquad
    I_{zz} = \int \left(x^{2} + y^{2}\right) \mathrm{d}m
    \label{eq:moments_of_inertia}
\end{equation}
are the moments of inertia about the $x_B$, $y_B$, and $z_B$ axes, respectively, and the off-diagonal entries
\begin{equation}
    I_{xy} = \int xy \,\mathrm{d}m, \qquad
    I_{xz} = \int xz \,\mathrm{d}m, \qquad
    I_{yz} = \int yz \,\mathrm{d}m
    \label{eq:products_of_inertia}
\end{equation}
are the corresponding products of inertia. All six quantities are taken about the centre of gravity and have units of $\mathrm{slug}\cdot\mathrm{ft}^{2}$. Since the aircraft is symmetric about its $x_B$--$z_B$ plane, the products of inertia $I_{xy}$ and $I_{yz}$ vanish, so that $I_{xz}$ is the only non-zero off-diagonal element of $\mathbf{J}$. As $\mathbf{J}$ is symmetric and positive definite, the quantity $\Gamma = I_{xx}I_{zz} - I_{xz}^{2}$ is strictly positive. With these definitions, the coefficients are expanded as
\begin{equation}
    \begin{aligned}
        & C_1 = \frac{(I_{yy} - I_{zz})I_{zz} - I_{xz}^2}{\Gamma}, \quad C_2 = \frac{(I_{xx} - I_{yy} + I_{zz})I_{xz}}{\Gamma}, \\  
        & \quad C_3 = \frac{I_{zz}}{\Gamma}, \quad C_4 = \frac{I_{xz}}{\Gamma}, \quad
        C_5 = \frac{I_{zz} - I_{xx}}{I_{yy}}, \quad C_6 = \frac{I_{xz}}{I_{yy}}, \\
        &\quad C_7 = \frac{1}{I_{yy}}, \quad C_8 = \frac{(I_{xx} - I_{yy})I_{xx} + I_{xz}^2}{\Gamma}, \quad C_9 = \frac{I_{xx}}{\Gamma}.
    \end{aligned}
    \label{eq:inertia_coeffs}
\end{equation}

For the F-16 model used in this work, these inertia quantities are updated by JSBSim \cite{jsbsim} at every integration step from the airframe inertia tensor together with the pilot and fuel point masses, and therefore depend on the instantaneous loading state. At the loading state used for the initial conditions of the engagements, that is, with an initial fuel load of $3000$ lb giving a total weight of $20{,}630$ lb and $m = 641.2$ slugs, the moments and the product of inertia, expressed in $\mathrm{slug}\cdot\mathrm{ft}^{2}$, are
\begin{equation}
    I_{xx} = 1.229\times10^{4}, \quad
    I_{yy} = 5.711\times10^{4}, \quad
    I_{zz} = 6.707\times10^{4}, \quad
    I_{xz} = 1.060\times10^{3}.
    \label{eq:f16_inertia}
\end{equation}
The values given in Eq.~(\ref{eq:f16_inertia}) vary with the fuel state; when the fuel is exhausted ($m = 548.0$ slugs, that is, $17{,}630$ lb), they become $I_{xx} = 9.50\times10^{3}$, $I_{yy} = 5.682\times10^{4}$, $I_{zz} = 6.411\times10^{4}$, and $I_{xz} = 9.46\times10^{2}\ \mathrm{slug}\cdot\mathrm{ft}^{2}$. The mass $m$ in Eq.~(\ref{eq:clmes}) is updated in the same manner.

JSBSim includes an inner-loop fly-by-wire (FBW) architecture. Unlike studies that directly output control-surface deflection commands, the continuous action vector output by the agent, $\mathbf{u}_{cmd} = [u_{roll}, u_{pitch}, u_{yaw}, u_{T}]^T$, represents pilot-level stick and throttle inputs that serve as target rate and G-load setpoints. Given the current flight state $\mathbf{x}$, the FBW system maps these reference signals into the actual actuator-level surface deflections and engine commands, denoted by $[\delta_a, \delta_e, \delta_r, \delta_T]^T$. This design is consistent with the human--machine interaction logic of a real cockpit and also allows the reinforcement learning agent to benefit implicitly from the built-in stability augmentation loops, thereby focusing on higher-level tactical decision-making rather than high-frequency disturbance rejection at the control-surface level.


\begin{figure}[h]
\centering
\includegraphics[width=0.6\textwidth]{graph/wez.png}
\caption{Definition of the weapons engagement zone.}\label{fig:wez}
\end{figure}

% \begin{figure*}[t]
%     \centering
%     \includegraphics[width=0.95\linewidth]{graph/gpu_env.png}
%     \caption{Multiprocessing parallel acceleration architecture.}
%     \label{fig:env}
% \end{figure*}

The weapon system and engagement criteria used in this work are inspired by the publicly available rules of the DARPA ADT. The environment defines a weapons engagement zone (WEZ) in front of the aircraft nose, shown in Figure~\ref{fig:wez}. Let the position of the ego aircraft be $\mathbf{p}_{agt}$ and that of the opponent aircraft be $\mathbf{p}_{oppo}$. The relative line-of-sight vector is defined as $\mathbf{d} = \mathbf{p}_{oppo} - \mathbf{p}_{agt}$, and the relative Euclidean distance is $D = \|\mathbf{d}\|_2$. If the unit nose-direction vector of the ego aircraft is denoted by $\mathbf{v}_{nose}$, then the track angle $\theta_{tra}$ is computed from the inner product as
\begin{equation}
    \theta_{tra} = \arccos \left( \frac{\mathbf{d} \cdot \mathbf{v}_{nose}}{\|\mathbf{d}\|_2} \right).
    \label{eq:tra}
\end{equation}
The opponent is considered to be inside the WEZ only when it lies within the effective firing range $[D_{\min}, D_{\max}]$ and inside a cone of half-aperture angle $\theta_{\max}$ around the gun line. This condition can be expressed by the indicator function $\mathbb{W}_{WEZ}(t)$:
\begin{equation}
    \mathbb{W}_{WEZ}(t) = 
    \begin{cases} 
        1, & \text{if } D_{\min} \le D(t) \le D_{\max} \text{ and } \theta_{tra}(t) \le \theta_{\max} \\ 
        0, & \text{otherwise} 
    \end{cases},
    \label{eq:wez}
\end{equation}
where $D_{\min}$ and $D_{\max}$ denote the near and far boundaries of the effective firing range and $\theta_{\max}$ denotes the maximum admissible off-boresight angle. These three thresholds are not free parameters of the learning problem, since they describe the geometry of a valid firing solution rather than the behaviour of the agent. They were determined by considering both the practical characteristics of within-visual-range gun engagements and the publicly released rules of the DARPA ADT. In this work they are set to $D_{\min} = 500$ ft, $D_{\max} = 4500$ ft and $\theta_{\max} = 15^\circ$. The near boundary corresponds to the minimum separation at which a tracking solution can still be established, and it rules out an unrealistic point-blank firing condition; the far boundary corresponds to the range beyond which the dispersion of the gun and the loss of lead accuracy make hits improbable at the closing speeds of the modelled aircraft; and the cone angle follows the ADT criterion, which requires a close alignment between the target and the gun line.

Both aircraft enter an engagement with an initial health of $\mathrm{HP}_{0} = 3$. When $\mathbb{W}_{WEZ}(t) = 1$, the agent applies continuous damage to the opponent at the current time step, and the decay of the opponent's health follows
\begin{equation}
    \text{HP}_{oppo}(t+\Delta t) = \max\left(0, \text{HP}_{oppo}(t) - \lambda(D) \cdot \mathbb{W}_{WEZ}(t) \cdot \Delta t \right)
    \label{eq:damage},
\end{equation}
where $\lambda(D) = \lambda_{0}\left(D_{\max}-D(t)\right)/\left(D_{\max}-D_{\min}\right)$ is the prescribed per-unit-time gun damage rate and $\lambda_{0} = 1$ HP/s is its value at the near boundary. The damage rate therefore decreases linearly with range and vanishes at the far boundary of the WEZ, so that a target held at the minimum range is destroyed within $3$ s, whereas a target engaged near the far boundary survives correspondingly longer for the same amount of tracking time.

At the same time, aerodynamic data, FBW computation, and combat logic are highly decoupled within the JSBSim framework, so the proposed simulation architecture has good extensibility. By replacing the aerodynamic database and model configuration, the environment can be extended to other aircraft types, such as the F-15, providing a basis for future studies on multi-aircraft and asymmetric air combat.

\subsection{Parallel Sampling Architecture}
\label{sec:parallel_acceleration}

In RL, environment sampling throughput is often the key bottleneck that limits training efficiency. From the perspective of the agent, the opponent policy should be treated as part of the environment transition. However, in practical implementation, directly coupling the opponent network into each simulation time step can substantially reduce training speed. Specifically, JSBSim flight dynamics simulation relies on serial floating-point computation on the CPU, whereas opponent policy inference is better suited to parallel tensor computation on the GPU. If CPU-based dynamics integration and single-sample neural network forward passes are executed alternately at every time step, the system suffers not only from frequent context switching, but also from the inability to exploit the batching advantage of the GPU. As a result, sampling efficiency is significantly degraded by computational serialization.

\begin{figure*}[t]
    \centering
    \includegraphics[width=0.95\linewidth]{graph/gpu_env.png}
    \caption{Multiprocessing parallel acceleration architecture.}
    \label{fig:env}
\end{figure*}

To address this issue, this work designs a multiprocessing parallel acceleration architecture based on heterogeneous-computing decoupling, whose core idea is to separate dynamics simulation from neural network inference. Figure~\ref{fig:env} shows the detail. First, the base environment is refactored into a pure JSBSim dynamics container with no embedded neural network modules. It exposes a state-action interface for both the ego aircraft and the opponent aircraft. Second, by combining CPU multiprocessing with GPU batched inference, a vectorized multi-process environment wrapper is constructed. Each subprocess independently executes flight dynamics simulation. At every interaction step the opponent observations from all sub-environments are collected, concatenated, and sent to the GPU for batched forward inference through the fixed opponent policy network. The resulting opponent actions are then distributed back to the corresponding CPU processes for the next simulation step. Finally, through an external interface abstraction, the wrapper hides the opponent observations and actions from the learning algorithm and exposes only the state and action interfaces of the ego agent. Therefore, under a fixed opponent policy, the overall training procedure remains externally equivalent to a standard single-agent Markov decision process.

This decoupled architecture effectively leverages heterogeneous computing resources, yielding a considerable improvement in simulation throughput on the same hardware platform. On the platform and experiment setup detailed in Section \ref{sec:experiment}, compared with the original serial implementation, whose maximum throughput is only $50$ simulation time steps per second, the proposed parallel acceleration structure increases the environment iteration speed to approximately $6700$ simulation time steps per second, yielding a $134\times$ speedup. This will provide computational support for subsequent reinforcement learning training.


\section{Methods}
\label{sec:RL_training}

PPO \cite{PPO} is a representative Actor--Critic policy optimization method and has been widely applied to decision-making tasks such as robotic control, flight control, and autonomous driving. The core idea of PPO lies in formulating a constrained policy update objective that prevents the new policy from deviating excessively from the old policy.

Let the current policy parameters be $\theta$ and the old policy parameters be $\theta_{\mathrm{old}}$. For a state $s_t$ and action $a_t$, the importance sampling ratio is defined as
\begin{equation}
r_t(\theta)=\frac{\pi_{\theta}(a_t|s_t)}{\pi_{\theta_{\mathrm{old}}}(a_t|s_t)}.
\end{equation}
PPO adopts a clipped surrogate objective function:
\begin{equation}
L^{\mathrm{clip}}(\theta)=\mathbb{E}_t\left[\min\left(r_t(\theta)\hat{A}_t,\ \mathrm{clip}(r_t(\theta),1-\epsilon,1+\epsilon)\hat{A}_t\right)\right],
\end{equation}
where $\hat{A}_t$ denotes the advantage function estimate and $\epsilon$ is the clipping coefficient.
By clipping the importance sampling ratio, this objective restricts the magnitude of the policy change in a single iteration, improving training stability. Meanwhile, a value function network is used to estimate state values and is jointly optimized with the policy network. Combining the policy term, the value term, and an entropy regularization term, the overall optimization objective can be written as
\begin{equation}
L(\theta,\phi)=L^{\mathrm{clip}}(\theta)-c_1L^{\mathrm{value}}(\phi)+c_2H(\pi_{\theta}),
\end{equation}
where $L^{\mathrm{value}}(\phi)$ represents the value function loss, $\phi$ is the value network parameters, $H(\pi_{\theta})$ is the policy entropy. The weight coefficients $c_1$ and $c_2$ are used to balance the contributions of the value-function loss and the entropy regularization, respectively. In practice, $c_1$ is chosen such that the value loss has a comparable scale to the policy objective and does not dominate the policy update, while $c_2$ is usually set to a relatively small value to encourage sufficient exploration without preventing policy convergence. These coefficients are typically selected empirically according to the training stability and final performance, and can be tuned through validation experiments. 

To estimate the advantage, we employ generalized advantage estimation (GAE), which takes the form
\begin{equation}
\hat{A}_t=\sum_{l=0}^{T-t-1}(\gamma\lambda)^l\delta_{t+l},
\qquad
\delta_t=r_t+\gamma V(s_{t+1})-V(s_t),
\end{equation}
where $V$ represents the value network, $\gamma$ is the discount factor and $\lambda$ is the advantage estimation parameter.

Compared with standard policy gradient methods, it constrains the policy update magnitude, which helps mitigate the instability caused by drastic policy oscillations during training. Therefore, PPO is well suited as the foundational algorithmic framework for air combat control policy learning.

\subsection{Baseline}

This work adopts PPO as the baseline reinforcement learning algorithm and implements it based on Stable-Baselines3 \cite{stable-baselines3}. 

We pre-construct a pool of several level-flight opponents. These opponents are all trained using PPO and learn stable flight policies under different reward functions. They are designed to maintain sustained flight under randomly initialized headings without losing control or crashing. These opponents possess only basic flight capability and cannot track targets. Compared with stationary targets or a scripted opponent, this design introduces more realistic flight-dynamics perturbations during baseline training and improves the effectiveness of policy learning.

These level-flight opponents are trained using different combinations of several basic reward terms, which are listed in Table~\ref{tab:baseline_opponent_rewards}.

During baseline training, the agent randomly samples opponents from this fixed opponent pool. In practice, the baseline model can acquire basic target-tracking and kill capabilities against opponents from the training set. Nevertheless, when evaluated against unseen opponents or in self-play tests, the baseline model exhibits weak generalization. In many cases, it fails to locate the target and initiate an attack against opponents outside the training set. In some situations, it even crashes. To address this issue, this work introduces a curriculum learning strategy and a dynamic opponent pool.



\subsection{Curriculum Learning}

Curriculum learning is commonly used to mitigate the training difficulty caused by sparse rewards. It constructs a task sequence that progresses from simple to complex. In this section, we propose a curriculum learning method composed of an objective curriculum and a dynamic opponent curriculum. The former employs two progressively designed reward functions to ensure that the agent first acquires stable flight control and then gradually develops air combat tactics. The latter adaptively adjusts opponent difficulty to provide progressively increasing game pressure for policy improvement. This design is intended to improve training stability and convergence efficiency.

\subsubsection{Objective Curriculum}
\label{sec:reward_curriculum}

We implement the objective curriculum by progressively adjusting the reward function. The objective curriculum is structured in two successive stages, where the agent first masters stable flight and then learns air combat skills. The first stage includes only reward terms related to flight quality and does not involve relative engagement geometry or firing-related indicators. The second stage introduces reward terms related to relative geometry and gun engagement, so that the agent can learn tactical positioning and evasive maneuvering.

In both stages the reward is composed of a small number of normalized terms, each of which represents one requirement of the task and is scaled by a weight, so that the contribution of the individual requirements can be analysed and adjusted separately. The general expressions of the terms of the two stages are given in Table~\ref{tab:reward_stage1} and Table~\ref{tab:reward_stage2}, respectively, together with the notation used for their thresholds and weights. The values assigned to these parameters, and the procedure by which they were calibrated, are described in Section~\ref{sec:reward_calibration}.


% \begin{table}
% \centering
% \caption{Stage-1 reward function.}
% \label{tab:reward_stage1}
% \begin{tabular}{@{}ll@{}}
% \toprule
% Reward Term & Mathematical Expression \\
% \midrule
% Aileron command penalty & 
% $r_{sa}=4\cdot\dfrac{0.5}{0.5+|\delta_a|}$ \\[6pt]
% Throttle command penalty & 
% $r_{st}=8\cdot\dfrac{0.1}{0.1+|\delta_t-0.4|}$ \\[6pt]
% Roll-angle penalty & 
% $r_{sr}=4\cdot\dfrac{0.5}{0.5+|\phi|}$ \\[6pt]
% Action smoothness penalty & 
% $r_{ap} = 4\left(\Phi_{ap}(s_{t+1})-\Phi_{ap}(s_t)\right)$, \\ 
% & $\Phi_{ap}(s)=-\underset{j}{\max}\,\dfrac{|\Delta u_j|}{u_{j,\max}}$ \\[6pt]
% Altitude stability penalty & 
% $r_{alt} = 8\times10^4\left(\Phi_{alt}(s_{t+1})-\Phi_{alt}(s_t)\right)$，\\
% & $\Phi_{alt}(s)=-\max\,\dfrac{|\Delta h|}{h_{\max}}$ \\
% \botrule
% \end{tabular}
% \end{table}
\begin{table}[h]
\caption{Stage-1 reward terms and their general expressions.}\label{tab:reward_stage1}%
\begin{tabular*}{\textwidth}{@{\extracolsep\fill}p{0.32\textwidth}p{0.64\textwidth}}
\toprule
Reward term & Mathematical expression \\
\midrule
Aileron command penalty
& $r_{sa}=w_{sa}\dfrac{s_{sa}}{s_{sa}+|\delta_a-\delta_a^{\ast}|}$ \\[6pt]

Throttle command penalty
& $r_{st}=w_{st}\dfrac{s_{st}}{s_{st}+|\delta_T-\delta_T^{\ast}|}$ \\[6pt]

Roll-angle penalty
& $r_{sr}=w_{sr}\dfrac{s_{sr}}{s_{sr}+|\phi-\phi^{\ast}|}$ \\[6pt]

Action smoothness penalty
& $r_{ap}=w_{ap}\left(\Psi_{ap}(t)-\Psi_{ap}(t-1)\right)$, 
  $\Psi_{ap}(t)=-\max_j\dfrac{\overline{|\Delta u_j|}_{L_{ap}}}{u_{j,ub}}$ \\[6pt]

Altitude stability penalty
& $r_{alt}=w_{alt}\left(\Psi_{alt}(t)-\Psi_{alt}(t-1)\right)$,
  $\Psi_{alt}(t)=-\dfrac{\overline{|\Delta h|}_{L_{alt}}}{h_{ub}}$ \\
\botrule
\end{tabular*}
\end{table}

The subscript ``ub'' in Table~\ref{tab:reward_stage1} denotes the upper bound of the corresponding state variable, that is, $u_{j,ub}=1$ for the normalized surface commands and $h_{ub}=85000$ ft for the altitude; $\delta_a^{\ast}$, $\delta_T^{\ast}$ and $\phi^{\ast}$ are the reference values of the aileron input, the throttle input and the roll angle, respectively; $w_{\bullet}$ and $s_{\bullet}$ denote the weight and the scaling factor of a term; and $\overline{|\Delta x|}_{L}$ denotes the mean absolute step-to-step change of the quantity $x$ over the preceding $L$ decision steps.

The first three terms are asymptotic error potentials. Each of them attains its maximum $w_{\bullet}$ when the corresponding variable equals its reference value and decays as the deviation grows, with $s_{\bullet}$ setting the deviation at which half of the maximum is retained. These terms therefore encode the requirement of a trimmed and level flight condition rather than a specific manoeuvre: the aileron and roll terms suppress unnecessary lateral control activity and residual bank, while the throttle term concentrates the engine thrust around the cruise setting, which prevents both the loss of airspeed that would reduce the available control authority and the sustained over-throttle condition that would accompany an aggressive energy strategy.

The two remaining terms are potential-difference terms, and they penalize the variation of the control activity and of the vertical motion rather than their absolute level. Since the potentials are defined as the windowed mean absolute rates of change, a steady manoeuvre leaves $\Psi$ almost unchanged and therefore produces almost no penalty, whereas rapid reversals of the control surfaces or abrupt changes of the vertical rate produce a penalty proportional to the corresponding increment. This formulation thus acts as a smoothness regularizer that does not discourage sustained manoeuvring, and, because the potentials depend only on the variation of the state, it does not bias the policy towards any particular engagement geometry. The weights of the two terms differ by more than four orders of magnitude, which reflects the different scales of the quantities involved: the normalized control deviation is of order $10^{-1}$ to $10^{0}$, whereas the normalized altitude increment of a typical manoeuvre is of order $10^{-4}$, so that a common weight would render the altitude term numerically inactive.

% \begin{table}
% \centering
% \caption{Stage-2 reward function.}
% \label{tab:reward_stage2}
% \begin{tabular}{@{}ll@{}}
% \toprule
% Reward Term & Mathematical Expression \\
% \midrule
% Relative position reward & 
% $\begin{aligned}
% r_{rp}&=\Bigl(\frac{\theta_{tra}}{\pi}-2\Bigr)\sigma\!\Bigl(\frac{\theta_{adv}}{\pi};18,0.5\Bigr)\\
% &\quad -\frac{\theta_{tra}}{\pi}+1
% \end{aligned}$ \\[10pt]
% Closure rate reward & 
% $\begin{aligned}
% r_{cr}&=\frac{V_{close}}{500}\Bigl[1-\sigma\!\Bigl(\frac{\theta_{adv}}{\pi};18,0.5\Bigr)\Bigr]\\
% &\quad \times\sigma\!\Bigl(d;\frac{1}{500},2900\Bigr)
% \end{aligned}$ \\[10pt]
% Aiming reward & 
% $r_{gb}=\Gamma_{B}(d)\Bigl[1-\sigma\!\Bigl(\frac{\theta_{tra}}{\pi};10^5,\frac{1}{180}\Bigr)\Bigr]$ \\[8pt]
% Opponent aiming penalty & 
% $r_{gr}=-\Gamma_{R}(d)\,\sigma\!\Bigl(\frac{\theta_{adv}}{\pi};800,\frac{178}{180}\Bigr)$ \\[8pt]
% Low-altitude penalty & 
% $r_{deck}=-4\Bigl[1-\sigma\!\Bigl(h;\frac{1}{20},1300\Bigr)\Bigr]$ \\
% \botrule
% \end{tabular}
% \end{table}
\begin{table}[h]
\caption{Stage-2 reward terms and their general expressions.}\label{tab:reward_stage2}%
\begin{tabular*}{\textwidth}{@{\extracolsep\fill}p{0.30\textwidth}p{0.66\textwidth}}
\toprule
Reward term & Mathematical expression \\
\midrule
Relative position reward
& $\begin{aligned}
r_{rp}
&=w_{rp}\left[\left(\hat{\theta}_{tra}-2\right)
\sigma\left(\hat{\theta}_{adv};k_{a},\hat{\theta}_{a}\right)
-\hat{\theta}_{tra}+1\right]
\end{aligned}$ \\[10pt]

Closure rate reward
& $\begin{aligned}
r_{cr}
&=w_{cr}\frac{V_{c}}{V_{ref}}
\left[1-\sigma\left(\hat{\theta}_{adv};k_{a},\hat{\theta}_{a}\right)\right]
\sigma\left(D;k_{D},D_{c}\right)
\end{aligned}$ \\[10pt]

Aiming reward
& $r_{gb}=\Gamma_{B}(D)\left[1-\sigma\left(\hat{\theta}_{tra};k_{g},\hat{\theta}_{g}\right)\right]$ \\[8pt]

Opponent aiming penalty
& $r_{gr}=-\Gamma_{R}(D)\,\sigma\left(\hat{\theta}_{adv};k_{r},\hat{\theta}_{r}\right)$ \\[8pt]

Low-altitude penalty
& $r_{deck}=-w_{deck}\left[1-\sigma\left(h;k_{h},h_{deck}\right)\right]$ \\
\botrule
\end{tabular*}
\end{table}

In Table~\ref{tab:reward_stage2}, $\sigma(x;k,x_0)=1/(1+e^{-k(x-x_0)})$ is the logistic function, in which $k$ controls the slope of the transition and $x_0$ its centre, and the normalized angles $\hat{\theta}_{tra}=\theta_{tra}/\pi$ and $\hat{\theta}_{adv}=\theta_{adv}/\pi$ are used so that all angular thresholds are expressed on a common scale. The adverse angle $\theta_{adv}$ measures the angular advantage of the ego aircraft: it is zero when the ego aircraft is directly behind the opponent, and it increases to $\pi$ when the opponent has reversed the situation and the ego aircraft lies in front of it. A small adverse angle therefore corresponds to an offensive situation, and a large one to a defensive situation.

The aiming reward and the opponent aiming penalty are weighted by the functions $\Gamma_{B}(D)$ and $\Gamma_{R}(D)$, which represent the range dependence of the firing opportunities of the two aircraft. Both are band-pass functions of range of the form
\begin{equation}
    \Gamma_{i}(D)=\beta_{i}\min\left\{\sigma\left(D;k_{i,1},D_{i,1}\right),\;1-\sigma\left(D;k_{i,2},D_{i,2}\right)\right\},
    \qquad i\in\{B,R\},
    \label{eq:range_weight},
\end{equation}
where the two logistic factors describe the near and the far boundary of the effective firing band and $\beta_{i}$ is the value of the function inside that band. A firing solution is thus treated as being most valuable in the intermediate range band, where a tracking solution can be established and sustained, and as being less valuable at very short range, where the geometry changes rapidly and an overshoot is imminent, and beyond the far boundary, where hits become improbable.

The relative position reward $r_{rp}$ is the primary shaping term for the angular advantage. It is bounded in $[-w_{rp},w_{rp}]$, and it attains $+w_{rp}$ when the ego aircraft is simultaneously behind the opponent ($\hat{\theta}_{adv}\to0$) and pointing its nose at it ($\hat{\theta}_{tra}\to0$), whereas it attains $-w_{rp}$ when the opponent has established the same advantage over the ego aircraft, because then $\sigma\to1$ and the term reduces to $-1$ independently of $\hat{\theta}_{tra}$. The logistic factor with slope $k_{a}$ and centre $\hat{\theta}_{a}$ separates the offensive and the defensive regime smoothly, so that the term provides a continuous gradient over the whole engagement instead of a binary assessment of the relative geometry.

The closure rate reward $r_{cr}$ complements the angular term by shaping the approach. It is proportional to the closure rate normalized by $V_{ref}$, so that reducing the distance is rewarded while the ego aircraft is in the offensive regime, and it is gated by two logistic factors: the first one removes the reward when the ego aircraft is not behind the opponent, which prevents the agent from being rewarded for closing while it is itself being pursued, and the second one attenuates the reward below the distance $D_{c}$, which discourages the agent from driving into an overshoot. The parameter $D_{c}$ therefore acts as a soft no-overshoot boundary at which the agent is expected to have established a firing position rather than to continue to reduce the distance.

The two aiming terms represent the firing opportunity of the two aircraft. Because the damage inflicted by the gun enters the return only through the outcome of the engagement, these terms provide the dense signal that guides the nose of the aircraft onto the firing line. The aiming reward $r_{gb}$ is active when the track angle falls below the alignment threshold $\hat{\theta}_{g}$, so that it rewards a settled gun solution rather than a mere pointing of the aircraft towards the opponent, and it is weighted by $\Gamma_{B}(D)$ so that the reward is concentrated in the range band where such a solution is sustainable. The opponent aiming penalty $r_{gr}$ mirrors this construction for the adversary: it becomes active when the adverse angle exceeds the threshold $\hat{\theta}_{r}$, that is, when the opponent's nose is aligned with the ego aircraft, and it is weighted by $\Gamma_{R}(D)$, so that it imposes an increasing cost on conceding the offensive position inside the opponent's firing envelope.

Finally, the low-altitude penalty $r_{deck}$ acts as a flight-safety constraint. It is negligible while the aircraft remains above $h_{deck}$ and grows rapidly below that altitude, so that it prevents the policy from trading an excessive amount of altitude for speed or from entering the ground-proximity regime during aggressive manoeuvring. The progressive design of the two stages, together with the terminal terms given below, enables the agent to develop tactically meaningful offensive and defensive behaviour. The structure of the stage-2 reward follows the approach adopted in the ADT project \cite{pope_hierarchical_2023}.

Besides the per-step terms of the two stages, the outcome of the engagement is accounted for by the terminal terms listed in Table~\ref{tab:reward_terminal}, which are applied once when an episode terminates.

\begin{table}[h]
\caption{Terminal reward terms, applied when an episode terminates.}\label{tab:reward_terminal}%
\begin{tabular*}{\textwidth}{@{\extracolsep\fill}p{0.25\textwidth}p{0.69\textwidth}}
\toprule
Terminal term & Mathematical expression \\
\midrule
Health margin
& $r_{hp}=w_{hp}\left(\mathrm{HP}(t_{f})-\mathrm{HP}_{oppo}(t_{f})\right)$ \\[6pt]

Kill bonus
& $r_{kill}=s_{kill}\,w_{kill}\,\tau_{rem}$ \\[6pt]

Loss penalty
& $r_{loss}=-\,s_{loss}\,w_{loss}\,\tau_{rem}$ \\
\botrule
\end{tabular*}
\end{table}

In Table~\ref{tab:reward_terminal}, $t_{f}$ denotes the terminal time of the episode, $\tau_{rem}$ the number of decision steps remaining before the maximum episode length is reached, and $s_{kill}$ and $s_{loss}$ are indicators that equal one when the opponent and the ego aircraft are destroyed in the air, respectively. The health margin is evaluated whenever neither aircraft has crashed, and it converts the residual health of the two aircraft into a continuous outcome signal, so that an engagement which reaches the time limit is not scored as a neutral event but reflects the damage exchanged during it. The kill bonus rewards the destruction of the opponent, and the loss penalty penalizes the destruction of the ego aircraft. Since both are proportional to the remaining time, an engagement that is decided early is rewarded, or penalized, more strongly than one that is decided late, which introduces the time pressure that is characteristic of a gun engagement: a firing solution that is obtained early leaves the opponent with less opportunity to reverse the situation, whereas a delayed solution consumes the time budget available for the engagement.

\subsubsection{Reward Composition and Parameter Calibration}
\label{sec:reward_calibration}

The parameters of the reward functions given in Tables~\ref{tab:reward_stage1}--\ref{tab:reward_terminal} fall into two groups, which are determined in different ways. The first group consists of the thresholds and scaling factors that describe the physical setting: the boundaries of the gun envelope in Eq.~(\ref{eq:wez}), the alignment thresholds and the range bands of the two aiming terms, the reference condition of trimmed flight, and the safety altitude. These quantities characterize the weapon, the flight envelope of the aircraft and the safety requirements of the manoeuvring airspace, and they are therefore fixed by the engagement scenario rather than being free parameters of the learning problem. The second group consists of the weights $w_{\bullet}$ and of the slopes $k_{\bullet}$ of the shaping terms. Since the terms interact, and since their numerical values are meaningful only relative to one another, these coefficients cannot be derived in closed form; they were calibrated iteratively, as described at the end of this subsection.

For the first stage, the reference values are $\delta_a^{\ast}=0$, $\phi^{\ast}=0$ and $\delta_T^{\ast}=0.4$, the last of which is the throttle setting prescribed at the initial condition of every episode. The scaling factors are $s_{sa}=s_{sr}=0.5$, so that a half-deflection aileron command or a bank angle of $0.5$ rad reduces the corresponding term to half of its maximum, and $s_{st}=0.1$, which concentrates the throttle term around the cruise setting and therefore penalizes large throttle excursions strongly. The weights are $w_{st}=8$ and $w_{sa}=w_{sr}=4$. The windows of the two potential-difference terms are $L_{ap}=10$ and $L_{alt}=20$ decision steps, that is, $1$ s and $2$ s of engagement time; a substantially shorter window would react to the numerical noise of the actuator commands, whereas a longer one would delay the response to genuine changes of the flight condition. Their weights are $w_{ap}=4$ and $w_{alt}=8\times10^{4}$, the latter compensating for the small normalized magnitude of the altitude increments, as discussed above.

For the second stage, the transition between the offensive and the defensive regime is centred at $\hat{\theta}_{a}=0.5$, that is, $\theta_{adv}=90^\circ$, with slope $k_{a}=18$, which spreads the transition over approximately $40^\circ$ and thereby avoids a discontinuous interpretation of the relative geometry. The closure rate is normalized by $V_{ref}=500$ ft/s, a value representative of the closing speeds encountered by the modelled aircraft in within-visual-range engagements, so that a typical approach produces a term of order unity; the no-overshoot attenuation is centred at $D_{c}=2900$ ft with slope $k_{D}=1/500$ ft$^{-1}$, so that the term is essentially removed inside the range where a tracking solution should already have been established. The alignment thresholds follow the geometry of the weapon: the aiming reward uses $\hat{\theta}_{g}=1/180$ together with $k_{g}=10^{5}$, so that it is active only when the nose is aligned with the firing line to within about $1^\circ$, whereas the opponent aiming penalty uses $\hat{\theta}_{r}=178/180$ with $k_{r}=800$, so that it grows as soon as the opponent's nose approaches the ego aircraft, with a gentler slope that provides a usable gradient before the opponent's firing solution is complete. The range weighting functions of Eq.~(\ref{eq:range_weight}) use $\beta_{B}=\beta_{R}=2.5$, $D_{B,1}=1000$ ft, $D_{B,2}=2900$ ft, $D_{R,1}=400$ ft and $D_{R,2}=4100$ ft, with slopes $k_{B,1}=k_{B,2}=1/50$ ft$^{-1}$, $k_{R,1}=1/35$ ft$^{-1}$ and $k_{R,2}=1/200$ ft$^{-1}$. The bands of the two functions thus occupy the middle of the corresponding firing ranges and retain a margin at both ends, which reflects the fact that a firing solution is transient at very short range and uncertain near the far boundary. The weight of the relative position term is $w_{rp}=1$, so that this term defines the unit of the tactical shaping scale, whereas the low-altitude penalty uses $w_{deck}=4$, that is, four times that unit, together with $h_{deck}=1300$ ft and $k_{h}=1/20$ ft$^{-1}$. The larger weight is deliberate, because entering this regime endangers the aircraft and therefore dominates any tactical consideration, while the threshold lies well below the altitude band in which the engagements are initialized ($10000\pm1000$ ft), so that the term remains inactive during normal manoeuvring. The terminal terms use $w_{hp}=10$, $w_{kill}=3.7$ and $w_{loss}=1$. With $\mathrm{HP}_{0}=3$, the health margin spans $\pm30$, whereas the kill bonus reaches approximately $2200$ when the opponent is destroyed at the beginning of an episode. The outcome of the engagement therefore dominates the accumulated shaping reward, which is essential in the present setting: because the gun inflicts damage continuously and the result of an engagement is decided by the residual health, a return dominated by per-step terms could be accumulated without ever resolving the engagement, a behaviour that the high proportion of drawn episodes reported in Section~\ref{sec:eval_benchmark} would otherwise encourage.

Because these coefficients interact, they were calibrated iteratively rather than optimized, by adjusting a small number of them at a time and inspecting the resulting behaviour both in the training curves and in the statistics of the evaluation episodes. The diagnostics used for this purpose are the per-term reward averages, which are recorded separately for every component during evaluation, together with the outcome statistics of the engagement, namely the win, draw and loss rates, the proportion of the engagement time during which the opponent remains inside the WEZ, the mean track angle, and the crash rate. The following observations determined the direction of the adjustments. When the recorded trajectories exhibited sustained oscillation of the control surfaces or an unsteady altitude trace, the weights of the corresponding rate terms, $w_{ap}$ and $w_{alt}$, were increased until the oscillations disappeared; conversely, when the agent became reluctant to manoeuvre and the engagements degenerated into long-range circling, these weights were reduced, because an excessive smoothness penalty suppresses the aggressive manoeuvring that the tactical terms are intended to encourage. When the policy retained a residual bank angle or an off-trim throttle setting, the weights of the corresponding asymptotic terms, $w_{sr}$, $w_{sa}$ and $w_{st}$, were increased. In the second stage, when the agent learned to point its nose at the opponent but failed to reduce the distance, the closure term was strengthened by increasing $w_{cr}$; when the recorded trajectories showed repeated overshoots at close range, that is, a rapid reversal of the closure rate after the distance had decreased below the gun range, $D_{c}$ was moved closer to the effective firing band and the weight of the aiming reward was increased, so that a settled firing solution was rewarded rather than the mere approach. When the agent spent a substantial part of the engagement in front of the opponent without attempting to disengage, the opponent aiming penalty was extended to longer ranges by increasing $D_{R,2}$ and its amplitude $\beta_{R}$ was increased, so that the cost of conceding the offensive position became comparable to the benefit provided by the angular terms. When the proportion of crashed episodes was not negligible, $h_{deck}$ was raised until the aircraft stopped entering the low-altitude regime. Finally, when the shaping terms indicated that the agent had learned to hold the offensive position while the proportion of drawn engagements remained high, $w_{kill}$ and $w_{hp}$ were increased, so that deciding the engagement replaced the accumulation of shaping reward as the dominant objective of the policy.

Each adjustment was validated by retraining the policy and re-evaluating it against a fixed set of opponents; an adjustment was retained only when it improved the outcome statistics without increasing the crash rate or degrading the quality of the flight, and it was reverted otherwise. The values reported above are those retained by this procedure, and they are listed together with the remaining experimental settings in Section~\ref{sec:experiment}. It should be noted that these values reflect the balance between the terms discussed in this section rather than a unique optimum, and they are reported in full so that the reward structure can be reproduced.

\subsubsection{Opponent Curriculum}
\label{sec:opponent_curriculum}

Objective curriculum cannot regulate the game difficulty during training. If the agent combats with an aggressive opponent from the very beginning, it may be shot down frequently before learning basic evasive maneuvers. As a result, the sparse reward makes it difficult for the policy network to extract useful gradients from a large number of failed episodes and potentially causing convergence to random behavior. To address this issue, we propose a dynamic-threshold opponent curriculum. By adjusting the opponent's target, the game difficulty increases smoothly.

% \begin{figure}
%     \centering
%     \includegraphics[width=0.49\linewidth]{graph/stage1.png}
%     \hfill
%     \includegraphics[width=0.49\linewidth]{graph/stage2.png}
%     \caption{Low-threat opponent (left) and high-threat opponent (right).}
%     \label{fig:oppo_curri}
% \end{figure}
\begin{figure}[h]
\centering
\includegraphics[width=0.9\textwidth]{graph/opponent_curriculum.png}
\caption{Low-threat opponent and high-threat opponent in the opponent curriculum.}\label{fig:oppo_curri}
\end{figure}

Opponent curriculum is achieved by changing how the opponent selects its target. Figure~\ref{fig:oppo_curri} shows a low-threat opponent and a high-threat opponent. Specifically, we define a threshold value $p(t)\in[0,1]$ that decays linearly with the training step $t$:
\begin{equation}
p(t) = \max\!\Bigl(0,\; 1-\frac{t}{T}\Bigr),
\end{equation}
where $T$ denotes the total number of steps over which the curriculum decays. During the initialization of each training episode, the system samples a random variable $\xi$ uniformly from $(0,1)$ and assigns a target point to the opponent according to the comparison between $\xi$ and the current threshold $p(t)$:
\begin{itemize}
\item If $\xi \leq p(t)$, the opponent is treated as a low-threat opponent. The system generates a target point randomly within the airspace, and the opponent follows a policy that flies toward this target point. In this mode, the opponent does not attack the ego and poses only a minimal threat.
\item If $\xi > p(t)$, the opponent is treated as a high-threat opponent. Its target point is updated in real time to the current position of the ego. In this case, the opponent behaves aggressively and continuously attempts to obtain gun-firing opportunities, imposing substantially greater adversarial pressure on the ego.
\end{itemize}

As training proceeds, $p(t)$ decreases linearly from 1 to 0, so the probability of encountering a high-threat opponent increases correspondingly. At the beginning of training, the agent faces almost exclusively harmless opponents flying toward random points, which allows it to safely explore the flight envelope and accumulate basic flight experience. As the training time increases, high-threat opponents begin to appear with increasing frequency, forcing the policy to gradually learn defensive, evasive, and counter-engagement behaviors. This difficulty scheduling design enables a smooth transition from non-adversarial flight to adversarial dogfighting. It helps alleviate the instability caused by introducing strong opponents too early, while also reducing the risk of policy stagnation induced by prolonged exposure to weak opponents.

\subsection{Dynamic Opponent Pool}
\label{sec:dynamic_opponent_pool}

The opponent curriculum described above regulates the training difficulty by changing the opponent's target assignment. However, this mechanism only changes the opponent's behavioral objective during an episode, while the opponent policy itself, i.e., its model parameters and intrinsic combat capability, remains unchanged. Therefore, although opponent curriculum can provide a smooth difficulty transition in the early stage of training, it is insufficient to continuously improve the adversarial competence of the opponent.

\begin{figure*}[t]
    \centering
    \includegraphics[width=0.95\linewidth]{graph/opponent_pool.png}
    \caption{Dynamic opponent pool.}
    \label{fig:opponent_pool}
\end{figure*}

To further enhance training, the ego policy should be exposed to opponents with progressively stronger and more diverse combat abilities. Training against a fixed opponent often causes the policy to overfit to a particular behavior pattern, whereas pure self-play may induce policy oscillation due to cyclic dominance. To address this issue, this study proposes a dynamic opponent pool. Fixed benchmark opponents are combined with competitive historical snapshots to provide the policy with a continuously evolving and behaviorally diverse set of adversaries. Details are shown in Figure~\ref{fig:opponent_pool}.

\subsubsection{Opponent Pool Initialization}

The initialization of the opponent pool is critical for stable policy learning. If the initial opponents are overly simple, the agent may overfit to limited behavior patterns, and when stronger opponents are introduced later, the learned policy and value estimates may transfer poorly, leading to performance degradation. Conversely, if the initial opponents are too difficult, the agent may fail to obtain sufficient meaningful interactions. To alleviate this cold-start problem, this paper introduces a goalpoint navigation policy as an initial protected opponent.

% At the early stage of training, the dynamic opponent pool has not yet been fully populated, and the policy interacts primarily with the protected opponents. If the initial protected opponents consist only of straight flight policies, the agent may overfit. When strong opponents are introduced later, the learned policy and value estimates often transfer poorly, which can lead to performance degradation. To alleviate this cold-start problem, this paper introduces a goalpoint navigation policy as an initial protected opponent.

The goalpoint model is defined as follows: from an arbitrary initial state, the aircraft will fly toward a randomly generated target point in space. Its observation consists of the ego aircraft state together with the geometric relationship between the aircraft and the target point. Instead of using the absolute coordinates of the target point, the model takes relative features as input, including relative azimuth, elevation, and distance. This design prevents the policy from memorizing specific spatial locations.

To facilitate the learning of this navigation task, a composite reward function is designed. 
Four general reward terms are used, namely an aileron magnitude penalty, an action smoothness penalty, a low-altitude penalty, and a closure-rate error term. The first two and the low-altitude penalty are of the same form as the corresponding terms of Tables~\ref{tab:reward_stage1} and~\ref{tab:reward_stage2}, whereas the closure-rate error term is specific to this auxiliary task: it is an asymptotic error potential whose target is a prescribed closure rate $V_{c}^{\ast} = 700$ ft/s, so that the aircraft approaches the target point at a moderate rate instead of accelerating towards it. The parameters of these four terms are $w_{ail}=1.0$ for the aileron magnitude penalty, $w_{ap}=3.5$ with $L_{ap}=10$ for the action smoothness penalty, $w_{deck}=4$, $h_{deck}=1300$ ft and $k_{h}=1/20$ ft$^{-1}$ for the low-altitude penalty, and $V_{c}^{\ast}=700$ ft/s with scaling factor $s=800$ ft/s and unit weight for the closure-rate error term.
In addition to these terms, the reward function incorporates two task-specific components to actively guide the aircraft: a distance reward and a heading reward. 
Both components are computed using a scaled asymptotic error form,
\begin{equation}
    r_i = w_i \left(1 - \frac{|x_i - x_i^\ast|/s_i}{|x_i - x_i^\ast|/s_i + 1}\right)
        = w_i \frac{s_i}{|x_i - x_i^\ast| + s_i},
    \label{eq:goalpoint_general}
\end{equation}
where $x_i$ is the current value of the corresponding variable, $x_i^\ast$ is the target value, $s_i$ is the scaling factor, and $w_i$ is the weight of the component. The form of Eq.~(\ref{eq:goalpoint_general}) is identical to that of the first three terms of Table~\ref{tab:reward_stage1}, so that the same interpretation applies: each component equals its weight $w_i$ when the corresponding error vanishes, retains half of that value when the error equals $s_i$, and decays smoothly for larger errors without introducing a singularity.

For the distance reward, the target distance is zero, i.e., $x_i^\ast=0$. Let $d(t)$ denote the relative distance between the aircraft and the target point in feet. Following Eq.~(\ref{eq:goalpoint_general}), the distance reward is defined as
\begin{equation}
    r_{\mathrm{dist}}(t) = w_{\mathrm{dist}}\frac{s_{\mathrm{dist}}}{d(t) + s_{\mathrm{dist}}},
    \label{eq:goalpoint_dist}
\end{equation}
where $w_{\mathrm{dist}} = 4.0$ and $s_{\mathrm{dist}} = 2000$ ft. The scaling factor is comparable to the separation at which the target points are generated, which are placed several thousand feet away from the aircraft in the horizontal plane and between $8000$ and $20000$ ft in altitude, so that the term retains an informative gradient over the whole approach and decays smoothly as the target point is reached; the large weight makes this term the dominant component of the reward of the auxiliary task, so that reaching the target point, rather than refining the flight condition, is its primary objective.

For the heading reward, since the desired heading error is zero, the reward is likewise given by
\begin{equation}
    r_{\mathrm{head}}(t) = w_{\mathrm{head}}\frac{s_{\mathrm{head}}}{|\theta_{tra}(t)| + s_{\mathrm{head}}},
    \label{eq:goalpoint_head}
\end{equation}
where $\theta_{tra}(t)$ is the angle between the nose of the aircraft and the line of sight to the target point, and the parameters are $w_{\mathrm{head}} = 0.2$ and $s_{\mathrm{head}} = 0.2$ rad, that is, approximately $11.5^\circ$. The scaling factor is comparable to the angular offset that the aircraft must remove in order to fly towards a target point, and the resulting term decreases from its maximum of $1$ when the aircraft points directly at the target point to approximately $0.1$ when the target point lies abeam, so that it acts as a persistent pointing signal which refines the direction of the flight without competing with the distance term, whose weight is twenty times larger. The target point is considered to be reached when the distance falls below $500$ ft, at which point the episode terminates.

In this work, the goalpoint model is pretrained separately using PPO, and the target point is randomly reset at every episode. The resulting goalpoint policy can guide the aircraft toward arbitrary target locations in the airspace using smooth control commands. Using this pretrained goalpoint model as the opponent, a batch of preliminary combat-capable models is trained by applying the reward function detailed in Table~\ref{tab:reward_stage2}. These policies are then evaluated according to their training curves, flight safety, and behavioral characteristics. Models with reliable training performance, low crash frequency, clear offensive or evasive behaviors, and diverse maneuvering styles are selected to form the initial protected opponent set. In this way, the initialized opponent pool contains not only the original goalpoint policy, but also several derived opponents that already possess basic maneuvering and combat interaction capability. Compared with initializing the pool using only handcrafted opponents, this procedure provides a broader opponent distribution. Furthermore, because the navigation targets of all opponents within the pool can be flexibly overridden, the system can seamlessly implement the opponent curriculum by toggling opponent behaviors between low threat and high threat. Combined with the objective curriculum, this integrated pipeline systematically guides the agent from basic flight control to complex air combat mastery.
% As a result, the dynamic opponent pool obtains richer adversarial diversity at the early stage of training, which provides a more effective initialization for the subsequent incorporation of historically strong policies.

\subsubsection{Opponent Pool Structure and Update}

Specifically, the proposed framework maintains an opponent pool with fixed capacity. At the beginning of each training episode, an opponent is independently sampled from the pool for each parallel environment. The opponent pool consists of two components. The first is a set of protected opponents initialized using the procedure described in the previous subsection, which provide stable baselines. The second is a set of dynamically updated policy slots used to store historically strong versions of the policy during training. The protected opponents cover representative behavior patterns with different levels of complexity and tactical styles. They remain unchanged throughout training so as to prevent the performance of the opponent pool from degrading due to frequent replacement. In contrast, the dynamic slots are filled progressively as training proceeds and are updated through periodic evaluation.

Let the evaluation interval be denoted by $K$. Whenever the number of training steps reaches an integer multiple of $K$, the current policy is frozen to form an evaluation snapshot, which is then matched against every opponent in the pool for $N$ independent episodes. To avoid interference from reward shaping, the update criterion does not rely directly on training return. Instead, it is constructed only from game outcomes and the remaining health at the end of engagement. For opponent $i$, the corresponding score is defined as
\begin{equation}
S_i = w_i - l_i + 0\cdot d_i + 0.5f_i + 0.01\bigl(hp_i - hp_i^{oppo}\bigr),
\end{equation}
where $w_i$, $l_i$, and $d_i$ denote the win rate, loss rate, and draw rate, respectively; $f_i$ denotes the opponent self-crash rate; and $hp_i$ and $hp_i^{oppo}$ denote the average remaining health percentages of the ego and the opponent aircraft, respectively. In this way, it offers a more direct measure of actual dogfight effectiveness and reduces the influence of reward hacking.

The update rule is as follows. If the opponent pool still contains empty dynamic slots, the current policy snapshot is inserted directly into one of them. If all dynamic slots are already occupied, the current policy is compared with the replaceable opponents in the pool, and it replaces the weakest one if its performance is superior. The protected opponents are never removed. 

By preserving fixed benchmark opponents together with a limited number of high-quality historical snapshots, the dynamic opponent pool helps balance adversarial diversity, training stability, and computational cost. Moreover, the design is compatible with the parallel sampling framework: different sub-environments independently sample opponents at each episode, so that a single batch of experience contains multiple adversarial styles. This design alleviates the bias caused by a single-opponent distribution and improves the generalization of the policy.


In summary, the baseline is trained solely against non-reactive, level-flight opponents that cannot perceive the ego aircraft, and it operates without curriculum learning or a dynamic opponent pool. To achieve better combat performance, our proposed method first constructs an initial opponent pool using the goalpoint model to alleviate the cold-start problem. Subsequently, the policy is trained by integrating the curriculum learning alongside the dynamic opponent pool to ensure stable learning and robust generalization. The performance of this method will be systematically compared against the baseline in Section~\ref{sec:experiment}.


\section{Experiments and Evaluation}
\label{sec:experiment}
\subsection{Experimental Setup}

This section describes the training framework and hardware platform in the experiments. All experiments are implemented based on Stable-Baselines3 and are trained using PPO. Except for the curriculum learning and the opponent selection, all other training configurations are kept identical across different methods to ensure a fair comparison.

The experimental comparison is organized from two perspectives. First, to evaluate the contribution of each proposed component, we start from a baseline and progressively consider the use of curriculum learning and the dynamic opponent pool. This yields four MLP-based models: the baseline model without curriculum learning or dynamic opponent pool (Baseline), the model with the dynamic opponent pool only (BD), the model with curriculum learning only (BC), and the model using both components (BDC-MLP). Second, to study the influence of the policy architecture, we further compare two implementations of the complete BDC method, namely an MLP-based version (BDC-MLP) and an LSTM-based \cite{LSTM} version (BDC-LSTM).

For the component ablation experiments, the MLP-based policy architecture is used by default to ensure that the compared models differ only in the use of curriculum learning and the dynamic opponent pool. For the MLP-based policy, the policy and value networks are parameterized by multilayer perceptrons and share a common feature extractor. For the LSTM-based policy, recurrent layers are further introduced to capture temporal dependencies in the observation sequence. The detailed hyperparameter settings for both architectures are summarized in Tables~\ref{tab:algo_net} and~\ref{tab:env_stage}.


Experiments were conducted on a workstation running Ubuntu 24.04.2 LTS, equipped with an Intel\textsuperscript{\textregistered} Core\texttrademark{} i9-14900KF CPU and a NVIDIA GeForce RTX\texttrademark{} 5070 Ti GPU. 


\subsection{Curriculum Learning Ablation}
To evaluate the effect of the curriculum learning on training stability and convergence performance, an ablation study is conducted. The experiments record the evolution of the mean evaluate return of three methods over $3.0\times10^8$ environment interaction steps. The corresponding training curves are shown in Figure~\ref{fig:training_curve}.

\begin{figure}[h]
    \centering
    \includegraphics[width=0.95\linewidth]{graph/reward_curve_plot.png}
    \caption{Episode return of the Baseline, BD, and BDC(MLP) methods. "Baseline" denotes standard PPO with a fixed opponent pool; "D" and "C" represent the Dynamic opponent pool and Curriculum learning components, respectively.}
    \label{fig:training_curve}
\end{figure}

From the overall trend, the BDC method outperforms the baseline. Although the baseline exhibits some improvement at the early stage of training, it begins to plateau after approximately $5\times10^7$ steps, with the mean return remaining around 500 for a long period. BDC method converges to approximately the 1000--1250 range, suggesting that curriculum learning and dynamic opponent pool helps improve policy performance.

A more direct view of the role of curriculum learning in training stability can be obtained by comparing the BDC method with BD. After removing curriculum learning, the model still benefits from the dynamic opponent pool, but the training curve shows pronounced fluctuations and experiences clear performance collapses around $4\times10^7$ and $1.3\times10^8$ steps. During the early stages of training, the agent faces three challenges: complex flight control, air-combat objectives, and an evolving opponent. Forcing the agent to address these simultaneously renders the optimization problem difficult. Ultimately, this may lead to unstable policy updates or even catastrophic forgetting.

By contrast, the BDC method gradually increases task difficulty through the curriculum. This process significantly alleviates the exploration burden at the beginning of training, allowing the model to achieve higher final returns while avoiding large performance oscillations. Although some fluctuations still remain in the later training stage as the dynamic opponent pool continues to evolve, the overall variance is clearly lower than BD method. These results demonstrate that the curriculum learning enhances the stability and robustness of the training process.


\subsection{Evaluation}

\subsubsection{Benchmark Pool Evaluation}
\label{sec:eval_benchmark}

To evaluate the generalization ability and tactical performance of the models, we construct a standardized benchmark opponent pool. This evaluation pool contains eight opponent policies covering diverse engagement styles, including aggressive, evasive, and random-maneuvering behaviors. These benchmark opponents are selected snapshots obtained from various training iterations and configurations. Their diverse tactical styles are induced by adjusting the relative weights of reward terms and varying the initial engagement geometries during their training. For instance, policies trained from advantageous initial positions tend to develop aggressive offensive tactics, whereas those initialized in disadvantageous positions learn more conservative or evasive behaviors. To provide a more intuitive characterization of the benchmark opponent pool, representative engagement scenarios of the eight opponents are illustrated in Figure~\ref{fig:benchmark_show}, where the red aircraft denotes the benchmark opponent. As shown in the figure, the selected opponents exhibit different tactical preferences, including spiral offensive maneuvers, weaving attacks, and evasive disengagement behaviors.

\begin{figure}[h]
\centering
\includegraphics[width=0.9\textwidth]{graph/benchmark_show.png}
\caption{Representative engagement scenarios of the eight benchmark opponents (the red aircraft).}\label{fig:benchmark_show}
\end{figure}

Five models are compared: the baseline model (Baseline), a model with curriculum learning only (BC), a model with the dynamic opponent pool only (BD), and two versions of the BDC method, namely the multilayer perceptron version (BDC-MLP) and the long short-term memory version (BDC-LSTM). During evaluation, each model is matched against all eight opponents, with 1000 independent Monte Carlo simulations conducted for each matchup. In these simulations, the Monte Carlo sampling is specifically applied to the initial states of the aircraft, ensuring that both the starting position and attitude are randomly initialized for every episode to thoroughly assess the model's robustness.

\begin{figure}[!htbp]
\centering
\includegraphics[width=0.9\textwidth]{graph/model_win_rates.png}
\caption{Win, draw, and loss rates of different models against the benchmark opponent pool over $8\times 1000$ simulated engagements per model.}\label{fig:win_rate}
\end{figure}
\begin{figure}[!htbp]
\centering
\includegraphics[width=0.9\textwidth]{graph/model_capabilities.png}
\caption{Comparison of key air-combat metrics, including gun opportunity ratio, mean specific energy, mean track angle, and total damage to the opponent.}\label{fig:combat_metrics}
\end{figure}

Figure~\ref{fig:win_rate} shows the win/draw/loss distributions on the benchmark opponent pool. The Baseline is almost unable to secure victories, achieving a win rate of only 0.42\%, with the vast majority of engagements ending in draws. This indicates that training only against fixed simple opponents may lead to policy overfitting. After introducing curriculum learning and the dynamic opponent pool separately, the win rates increase to 22.4\% and 33.3\%, respectively, showing that both components can independently improve generalization performance. The BDC method further improves the overall adversarial performance, with BDC-MLP and BDC-LSTM achieving win rates of 41.0\% and 55.6\%, respectively. The BDC-LSTM model therefore outperforms the baseline by approximately 55.2 percentage points.

Fig.~\ref{fig:combat_metrics} further compares four key air-combat metrics. The gun opportunity ratio is defined as the proportion of the total engagement time during which the opponent remains within the aircraft’s weapons engagement zone. The mean specific energy denotes the aircraft's average specific energy (calculated as $E_s = h + \frac{V_T^2}{2g}$, where $h$ is the altitude, $V_T$ is the true velocity, and $g$ is the gravitational acceleration) over the course of the engagement. Maintaining a higher level of specific energy is important in dogfight scenarios \cite{energy1, energy2}, as it reflects the policy's ability to manage the conversion between kinetic and potential energy, which is often critical for gaining positional and tactical advantages during engagement. Mean track angle refers to the average track angle of the aircraft, where a smaller value indicates stronger target-following capability. The total damage to opponent denotes the average total damage to opponent in the whole engagement.

First, in terms of gun opportunity ratio, BDC-MLP and BDC-LSTM reach 16.0\% and 19.1\%, respectively, which are higher than 0.92\% achieved by the baseline. This indicates that the BDC method can maintain the opponent within an effective firing envelope more consistently. Second, for mean specific energy, the BDC method outperforms the compared methods, with BDC-LSTM reaching 22401~ft, suggesting better energy retention during high-maneuver engagements. Third, in mean track angle, BDC-LSTM obtains the best result (0.696~rad), indicating that it can maintain a more favorable target-pointing condition in complex engagements. Lastly, the BDC method also achieves the highest score in the total damage to opponent metric, inflicting the most damage on the adversary.


\subsubsection{Cross-Play Evaluation}
\label{sec:cross_play}

In addition to the benchmark opponent pool evaluation, we further conduct cross-play experiments to compare the relative combat performance of different models. The models are evaluated in a round-robin manner, with each model competing against every other model. We evaluated the win rate of each model and the difference in remaining HP between the ego and the opponent at the end of each engagement. The results are shown in Figure~\ref{fig:heatmap_crossplay}.

\begin{figure}[h]
    \centering
    \includegraphics[width=0.95\linewidth]{graph/crossplay_heatmap.png}
    \caption{Cross-play results among the ablation models. Each cell reports the win rate and the remaining-HP difference ($\Delta$HP) of the row agent against the column agent.}
    \label{fig:heatmap_crossplay}
\end{figure}

In Figure~\ref{fig:heatmap_crossplay}, the color indicates the HP difference between the row model and the column model. The results show that models using BDC method achieve higher win rates and positive $\Delta$HP values against all ablation models, whereas the baseline is consistently at a disadvantage when facing the other methods. Introducing curriculum learning alone or the dynamic opponent pool alone can both improve policy quality, but neither matches the BDC method. Furthermore, a comparison across network architectures reveals that LSTM-based policies outperform their MLP counterparts. %This indicates that curriculum learning and the dynamic opponent pool play clearly complementary roles, and their joint use further enhances practical combat performance.

\begin{figure}[h]
    \centering
    \includegraphics[width=0.95\linewidth]{graph/violin_combined.png}
    % \hfill
    % \includegraphics[width=0.48\linewidth]{graph/violin_energy.png}
    \caption{Distributions of remaining health points (left) and specific energy (right) in direct cross-play between BDC-LSTM and BDC-MLP. Violin width indicates data density. The embedded boxplot shows the median and interquartile range (IQR), with whiskers extending to $1.5\times\mathrm{IQR}$. Overlaid scatter points represent individual samples. Together, these elements visualize the distribution shape, dispersion, and sample-level variability of both metrics.}
    \label{fig:violin_hp_energy}
\end{figure}

% \begin{figure*}
%     \centering
%     % 第一行：1 和 4
%     \begin{subfigure}{0.45\textwidth}
%         \includegraphics[width=\linewidth]{graph/RS_finish.png}
%         \caption{Rolling scissors.}
%         \label{fig:traj_scissors}
%     \end{subfigure}
%     \hfill
%     \begin{subfigure}{0.48\textwidth}
%         \includegraphics[width=\linewidth]{graph/defensive_finish.png}
%         \caption{Defensive barrel roll.}
%         \label{fig:traj_barrel_roll}
%     \end{subfigure}
    
%     \vspace{0.3cm}
    
%     % 第二行：2 和 3
%     \begin{subfigure}{0.45\textwidth}
%         \includegraphics[width=\linewidth]{graph/low_yoyo_finish.png}
%         \caption{Low Yo-Yo.}
%         \label{fig:traj_low_yoyo}
%     \end{subfigure}
%     \hfill
%     \begin{subfigure}{0.45\textwidth}
%         \includegraphics[width=\linewidth]{graph/high_yoyo_finish.png}
%         \caption{High Yo-Yo.}
%         \label{fig:traj_high_yoyo}
%     \end{subfigure}
    
%     \caption{Three-dimensional trajectory visualizations with six-DOF aircraft attitudes during direct engagements between BDC-LSTM (blue) and BDC-MLP (red). Solid dots indicate the starting positions. The four subfigures show representative Basic Fighter Maneuvers (BFM) observed during the engagements, including rolling scissors, defensive barrel roll, low Yo-Yo, and high Yo-Yo.}
%     \label{fig:3d_trajectories}
% \end{figure*}
\begin{figure*}[t]
\centering
\includegraphics[width=0.95\textwidth]{graph/engagement.png}
\caption{Three-dimensional trajectory visualizations with 6-DOF aircraft attitudes during direct engagements between BDC-LSTM (blue) and BDC-MLP (red). Solid dots indicate the starting positions. The four panels show representative basic fighter maneuvers observed during the engagements: (a) rolling scissors, (b) defensive barrel roll, (c) low Yo-Yo, and (d) high Yo-Yo.}\label{fig:3d_trajectories}
\end{figure*}

We can also present the detailed engagement results between BDC-LSTM and BDC-MLP. Since both are trained with the proposed method and therefore have similar performance, the proportion of draws in their direct encounters is high. To further analyze their differences, Figure~\ref{fig:violin_hp_energy} presents the distributions of remaining HP and specific energy in cross-play.


As shown in Figure~\ref{fig:violin_hp_energy}, BDC-LSTM outperforms BDC-MLP in direct competition. Their win rates are 6.2\% and 0.9\%, respectively, and BDC-LSTM also attains a higher average remaining HP (2.6 vs.\ 2.3). Meanwhile, the average specific energy of BDC-LSTM is approximately 20388~ft, which is higher than the 13499~ft of BDC-MLP, indicating better energy management during engagements.

In addition, the two models exhibit different energy distribution patterns. BDC-MLP shows a concentrated unimodal distribution, whereas BDC-LSTM presents a bimodal characteristic. This suggests that the LSTM architecture may help the policy develop more differentiated energy-utilization modes under different engagement situations. These findings suggest that incorporating temporal memory may improve decision-making flexibility in adversarial combat.

% \subsubsection{Human-versus-agent Experiment}
% \label{sec:human_machine}

% To further evaluate the model's generalization to out-of-distribution opponents, human-versus-agent experiments were conducted. More than ten volunteers were invited to engage in real-time one-on-one WVR gun dogfights against the BDC-LSTM agent using joystick controllers. The purpose of this experiment was not to compare AI performance with that of professional pilots, but rather to examine the stability and robustness of the policy when facing unseen control styles and non-scripted strategies. Human control behavior varies substantially in reaction time, control habits, and maneuver selection, making it a representative out-of-distribution test setting.

% \begin{figure}
%     \centering
%     \includegraphics[width=0.85\linewidth]{graph/human_vs_agent.png}
%     \caption{Distribution of average remaining Health Points and mean Track Angle in real-time dogfights between human pilots and the BDC-LSTM agent.}
%     \label{fig:human_vs_agent}
% \end{figure}

% Figure \ref{fig:human_vs_agent} presents the results in the human-versus-agent evaluation. First, in terms of mean track angle, BDC-LSTM achieved lower track angle than the human participants, indicating that it was able to maintain stable target pointing and trajectory tracking against irregular maneuvers. Second, in terms of average remaining HP, the agent achieved a value of 2.941, compared with 2.293 for the human participants, suggesting that the agent maintained an advantage during the engagements.

% Overall, BDC-LSTM maintained stable performance when facing human opponents in this out-of-distribution test setting. These results suggest that the proposed training framework achieves a reasonable degree of generalization and robustness.


\subsection{Engagement Demonstration}
\label{sec:tactical_analysis}
To better analyze the agent's maneuvering decisions, we developed a three-dimensional trajectory visualization with six-degree-of-freedom state information (3D attitude plot). This section presents four representative engagement patterns observed in direct matches between BDC-LSTM (blue) and BDC-MLP (red).


Figure \ref{fig:3d_trajectories} illustrates four representative basic fighter maneuvers:
\begin{enumerate}
    \item \textbf{Rolling Scissors}: As shown in Figure~\ref{fig:3d_trajectories}(a), after a close-range merge, both aircraft attempt to gain a favorable firing position behind the opponent, and the engagement develops into a typical rolling scissors maneuver.
    
    \item \textbf{Defensive Barrel Roll}: As shown in Figure~\ref{fig:3d_trajectories}(b), the blue aircraft performs a defensive barrel roll while being pursued. By following a three-dimensional helical trajectory, it reduces the closure rate and changes the relative geometry, thereby inducing an overshoot by the red aircraft and improving its own tactical situation.
        
    \item \textbf{Low Yo-Yo}: As shown in Figure~\ref{fig:3d_trajectories}(c), when the pursuit distance begins to increase, the blue aircraft executes a dive to gain speed, shortens its path on a lower turning plane, and compresses the relative distance before pulling up again to recover an attacking posture.
    
    \item \textbf{High Yo-Yo}: As shown in Figure~\ref{fig:3d_trajectories}(d), when the blue aircraft enters the pursuit with a high closure rate, it does not simply follow the red aircraft in a planar turn. Instead, it deliberately climbs, converting part of its kinetic energy into potential energy, reducing the risk of overshoot and improving the turning geometry before re-entering a favorable firing position.
\end{enumerate}

These representative cases suggest that the BDC method is able to spontaneously develop diverse maneuvering strategies according to relative position, speed, and energy state. 

\section{Conclusion}
\label{sec:conclusion}

This paper proposes a reinforcement learning framework for within-visual-range autonomous air combat, providing a solution with generalization and robustness. Through the integration of curriculum learning and dynamic opponent pool, the proposed framework stabilizes training and mitigates cyclic dominance. Experimental evaluations indicate that the proposed method yields competitive performance against baseline methods in terms of win rate and energy management. Ultimately, this work offers a viable baseline for addressing flight control and tactical decision-making in autonomous dogfighting.


\begin{appendices}

\section{}

The appendix provides supplementary implementation details for the proposed training framework. Table~\ref{tab:full_state_space} summarizes the complete state observation space together with their corresponding physical lower and upper bounds. Table~\ref{tab:baseline_opponent_rewards} lists the basic reward terms used for training the level-flight opponents, where $\sigma(\cdot)$ denotes the logistic function; $u_{\mathrm{ail}}^\ast$, $u_{\mathrm{thr}}^\ast$, and $\phi^\ast$ denote the target values of aileron input, throttle input, and roll angle, respectively; $h_{\mathrm{low}}$ and $h_{\mathrm{high}}$ denote the lower and upper boundaries of the allowable flight envelope; $k_a,k_t,k_r,k_l,k_h,k_s,k_{\Delta h}$ are the weighting coefficients of the corresponding reward terms; and $s_l,s_h$ control the slope of the boundary penalty transition. Table~\ref{tab:algo_net} reports the main algorithm and network hyperparameters used for both the LSTM-based and MLP-based policies. Table~\ref{tab:env_stage} presents the environment and stage configuration.

\begin{table*}[h]
\caption{Complete state observation space variables and extremum boundaries.}\label{tab:full_state_space}%
\small
\begin{tabular*}{\textwidth}{@{\extracolsep\fill}p{0.25\textwidth}p{0.35\textwidth}cc}
\toprule
State category & Variable name & Lower bound & Upper bound \\
\midrule
Aircraft base states
& $h$ & $-1400$ ft & $85000$ ft \\
& $\theta$ & $-\pi/2$ rad & $\pi/2$ rad \\
& $\phi$ & $-\pi$ rad & $\pi$ rad \\
& $u$ & $-2200$ ft/s & $2200$ ft/s \\
& $v$ & $-2200$ ft/s & $2200$ ft/s \\
& $w$ & $-2200$ ft/s & $2200$ ft/s \\
& $p$ & $-2\pi$ rad/s & $2\pi$ rad/s \\
& $q$ & $-2\pi$ rad/s & $2\pi$ rad/s \\
& $r$ & $-2\pi$ rad/s & $2\pi$ rad/s \\
\midrule
Aircraft tracking states
& heading & $0$ deg & $360$ deg \\
& $a_x,a_y,a_z$ & -- & -- \\
& $a_{roll},a_{pitch},a_{yaw}$ & -- & -- \\
& $\alpha$ & $-180$ deg & $180$ deg \\
& $\beta$ & $-180$ deg & $180$ deg \\
& aileron, elevator, rudder & $-1$ & $1$ \\
& throttle Aug, engine thrust & -- & -- \\
& vtrue & $0$ ft/s & $2200$ ft/s \\
\midrule
Relative geometry
& $D$ & $0$ ft & $500000$ ft \\
& $\theta_{tra}$ & $-\pi$ rad & $\pi$ rad \\
& $\theta_{bea}$, accounting roll/pitch & $-\pi$ rad & $\pi$ rad \\
& $\theta_{ele}$, accounting roll/pitch & $-\pi$ rad & $\pi$ rad \\
& $\theta_{bea}$, point mass & $-\pi$ rad & $\pi$ rad \\
& $\theta_{ele}$, point mass & $-\pi$ rad & $\pi$ rad \\
& $\theta_{adv}$ & $-\pi$ rad & $\pi$ rad \\
& $V_c$ & $-100000$ ft/s & $100000$ ft/s \\
\midrule
Opponent states
& $x_{oppo},y_{oppo}$ & $-100000$ ft & $100000$ ft \\
& $h_{oppo}$ & $-1400$ ft & $85000$ ft \\
& $\phi_{oppo},\theta_{oppo}$ & $-\pi,-\pi/2$ rad & $\pi,\pi/2$ rad \\
& $\mathrm{heading}_{oppo}$ & $0$ deg & $360$ deg \\
& $u_{oppo},v_{oppo},w_{oppo}$ & $-2200$ ft/s & $2200$ ft/s \\
& $p_{oppo},q_{oppo},r_{oppo}$ & $-2\pi$ rad/s & $2\pi$ rad/s \\
& $\alpha_{oppo},\beta_{oppo}$ & $-180$ deg & $180$ deg \\
& $\theta_{tra,oppo}$ & $-\pi$ rad & $\pi$ rad \\
& $\mathrm{vtrue}_{oppo}$ & $0$ ft/s & $2200$ ft/s \\
& $\theta_{bea,oppo},\theta_{ele,oppo}$ & $-\pi$ rad & $\pi$ rad \\
\midrule
Historical actions
& $\delta_a,\delta_e,\delta_r,\delta_T$ & $-1$ & $1$ \\
\botrule
\end{tabular*}
\end{table*}


\begin{table}[h]
\caption{Basic reward terms used for training the level-flight opponents.}\label{tab:baseline_opponent_rewards}%
\begin{tabular*}{\textwidth}{@{\extracolsep\fill}p{0.32\textwidth}p{0.64\textwidth}}
\toprule
Reward term & Mathematical form \\
\midrule
Aileron magnitude penalty
& $r_{\mathrm{ail}}=-k_a |u_{\mathrm{ail}}-u_{\mathrm{ail}}^\ast|$ \\[4pt]

Throttle penalty
& $r_{\mathrm{thr}}=-k_t |u_{\mathrm{thr}}-u_{\mathrm{thr}}^\ast|$ \\[4pt]

Roll attitude penalty
& $r_{\mathrm{roll}}=-k_r |\phi-\phi^\ast|$ \\[4pt]

Lower altitude penalty
& $r_{\mathrm{low}}=-k_l\left(1-\sigma\left(\dfrac{h-h_{\mathrm{low}}}{s_l}\right)\right)$ \\[4pt]

Upper altitude penalty
& $r_{\mathrm{high}}=-k_h\left(1-\sigma\left(\dfrac{h_{\mathrm{high}}-h}{s_h}\right)\right)$ \\[4pt]

Action smoothness penalty
& $r_{\mathrm{smooth}}=-k_s \sum_{i\in\{\mathrm{ail},\mathrm{ele},\mathrm{rud}\}}
|u_i^{(t)}-u_i^{(t-1)}|$ \\[4pt]

Altitude smoothness penalty
& $r_{\Delta h}=-k_{\Delta h}|h_t-h_{t-1}|$ \\
\botrule
\end{tabular*}
\end{table}


\begin{table}
  \centering
  \caption{Algorithm and network hyperparameters.}
  \label{tab:algo_net}
  \begin{tabular}{ccc}
    \toprule
    Hyperparameter & LSTM & MLP \\
    \midrule
    Learning rate & $5.0\times10^{-5}$ & $1.0\times10^{-4}$ \\
    Rollout steps & 2048 & 2048 \\
    Batch size & 128 & 128 \\
    Epochs & 10 & 10 \\
    Gamma ($\gamma$) & 0.99 & 0.99 \\
    GAE lambda ($\lambda$) & 0.95 & 0.95 \\
    Clip range & 0.2 & 0.2 \\
    Value function coefficient & 0.5 & 0.5 \\
    Max grad norm & 0.5 & 0.5 \\
    Share features extractor & False & True \\
    LSTM hidden size & 256 & -- \\
    Network arch (policy) & [512, 512] & [12288] \\
    Network arch (value) & [512, 512] & Shared \\
    \botrule
  \end{tabular}
\end{table}

% \begin{table}
%   \centering
%   \caption{Environment and stage configuration.}
%   \label{tab:env_stage}
%   \begin{tabular}{cc}
%     \toprule
%     Parameter & Value \\
%     \midrule
%     Plane & F16 \\
%     Clip obs & 10.0 \\
%     Seed & 42 \\
%     Num CPUs & 30 \\
%     Eval frequency & 50,000 \\
%     Battle step & 10,000,000 \\
%     Battle cycles & 30 \\
%     Number of eval episodes & 50 \\
%     Threshold timesteps & 20,000,000 \\
%     Update frequency (episodes) & 1,000 \\
%     \botrule
%   \end{tabular}
% \end{table}
\begin{table}[h]
\caption{Environment and stage configuration.}\label{tab:env_stage}%
\begin{tabular}{@{}ll@{}}
\toprule
Parameter & Value \\
\midrule
Aircraft model & F-16 \\
Observation clipping threshold & 10.0 \\
Random seed & 42 \\
Number of CPUs & 30 \\
Evaluation frequency & 50,000 episodes \\
Battle step & 10,000,000 \\
Battle cycles & 30 \\
Number of evaluation episodes & 50 \\
Threshold timesteps & 20,000,000 \\
Update frequency & 1,000 episodes \\
\botrule
\end{tabular}
\end{table}

% \begin{table}
% \centering
% \caption{Basic reward terms used for training the level-flight opponents.}
% \label{tab:baseline_opponent_rewards}
% \begin{tabular}{ll}
% \toprule
% \makecell[tl]{Reward term} & Mathematical form \\
% \midrule
% \makecell[tl]{Aileron magnitude \\ penalty} &
% $r_{\mathrm{ail}}=-k_a |u_{\mathrm{ail}}-u_{\mathrm{ail}}^\ast|$ \\

% \makecell[tl]{Throttle \\ penalty} &
% $r_{\mathrm{thr}}=-k_t |u_{\mathrm{thr}}-u_{\mathrm{thr}}^\ast|$ \\

% \makecell[tl]{Roll attitude \\ penalty} &
% $r_{\mathrm{roll}}=-k_r |\phi-\phi^\ast|$ \\

% \makecell[tl]{Lower altitude \\ boundary penalty} &
% $r_{\mathrm{low}}=-k_l\!\left(1-\sigma\!\left(\dfrac{h-h_{\mathrm{low}}}{s_l}\right)\right)$ \\

% \makecell[tl]{Upper altitude \\ boundary penalty} &
% $r_{\mathrm{high}}=-k_h\!\left(1-\sigma\!\left(\dfrac{h_{\mathrm{high}}-h}{s_h}\right)\right)$ \\

% \makecell[tl]{Action smoothness \\ penalty} &
% $r_{\mathrm{smooth}}=-k_s \sum_{i\in\{\mathrm{ail},\mathrm{ele},\mathrm{rud}\}}
% |u_i^{(t)}-u_i^{(t-1)}|$ \\

% \makecell[tl]{Altitude smoothness \\ penalty} &
% $r_{\Delta h}=-k_{\Delta h}|h_t-h_{t-1}|$ \\
% \botrule
% \end{tabular}
% \end{table}

% In Table \ref{tab:baseline_opponent_rewards}, $\sigma(\cdot)$ denotes the logistic function; $u_{\mathrm{ail}}^\ast$, $u_{\mathrm{thr}}^\ast$, and $\phi^\ast$ denote the target values of aileron input, throttle input, and roll angle, respectively; $h_{\mathrm{low}}$ and $h_{\mathrm{high}}$ denote the lower and upper boundaries of the allowable flight envelope; $k_a,k_t,k_r,k_l,k_h,k_s,k_{\Delta h}$ are the weighting coefficients of the corresponding reward terms; and $s_l,s_h$ control the slope of the boundary penalty transition.


% Table \ref{tab:full_state_space} gives the state space definition and physical boundary.

% \begin{table*}
%     \centering
%     \caption{Complete state observation space variables and extremum boundaries.}
%     \label{tab:full_state_space}
%     \small % Appendix table font size can be reduced appropriately
%     \begin{tabular}{llcc}
%         \toprule
%         \textbf{State Category} & \textbf{Variable Name} & \textbf{Lower Bound} & \textbf{Upper Bound} \\
%         \midrule
%         \multirow{9}{*}{\textbf{Aircraft Base States}}
%         & $h$ & $-1400$ ft & $85000$ ft \\
%         & $\theta$ & $-\pi/2$ rad & $\pi/2$ rad \\
%         & $\phi$ & $-\pi$ rad & $\pi$ rad \\
%         & $u$ & $-2200$ ft/s & $2200$ ft/s \\
%         & $v$ & $-2200$ ft/s & $2200$ ft/s \\
%         & $w$ & $-2200$ ft/s & $2200$ ft/s \\
%         & $p$ & $-2\pi$ rad/s & $2\pi$ rad/s \\
%         & $q$ & $-2\pi$ rad/s & $2\pi$ rad/s \\
%         & $r$ & $-2\pi$ rad/s & $2\pi$ rad/s \\
%         \midrule
%         \multirow{8}{*}{\textbf{Aircraft Tracking States}}
%         & heading & $0$ deg & $360$ deg \\
%         & $a_x, a_y, a_z$ & -- & -- \\
%         & $a_{roll}, a_{pitch}, a_{yaw}$ & -- & -- \\
%         & $\alpha$ & $-180$ deg & $180$ deg \\
%         & $\beta$ & $-180$ deg & $180$ deg \\
%         & aileron, elevator, rudder & $-1$ & $1$ \\
%         & throttle\_Aug, engine\_thrust & -- & -- \\
%         & vtrue & $0$ ft/s & $2200$ ft/s \\
%         \midrule
%         \multirow{8}{*}{\textbf{Relative Geometry}}
%         & $D$ & $0$ ft & $500000$ ft \\
%         & $\theta_{tra}$ & $-\pi$ rad & $\pi$ rad \\
%         & $\theta_{bea}$ (accounting roll/pitch) & $-\pi$ rad & $\pi$ rad \\
%         & $\theta_{ele}$ (accounting roll/pitch) & $-\pi$ rad & $\pi$ rad \\
%         & $\theta_{bea}$ (point mass) & $-\pi$ rad & $\pi$ rad \\
%         & $\theta_{ele}$ (point mass) & $-\pi$ rad & $\pi$ rad \\
%         & $\theta_{adv}$ & $-\pi$ rad & $\pi$ rad \\
%         & $V_c$ & $-100000$ ft/s & $100000$ ft/s \\
%         \midrule
%         \multirow{10}{*}{\textbf{Opponent States}}
%         & $x_{oppo}, y_{oppo}$ & $-100000$ ft & $100000$ ft \\
%         & $h_{oppo}$ & $-1400$ ft & $85000$ ft \\
%         & $\phi_{oppo}, \theta_{oppo}$ & $-\pi, -\pi/2$ rad & $\pi, \pi/2$ rad \\
%         & $\text{heading}_{oppo}$ & $0$ deg & $360$ deg \\
%         & $u_{oppo}, v_{oppo}, w_{oppo}$ & $-2200$ ft/s & $2200$ ft/s \\
%         & $p_{oppo}, q_{oppo}, r_{oppo}$ & $-2\pi$ rad/s & $2\pi$ rad/s \\
%         & $\alpha_{oppo}, \beta_{oppo}$ & $-180$ deg & $180$ deg \\
%         & $\theta_{tra, oppo}$ & $-\pi$ rad & $\pi$ rad \\
%         & $\text{vtrue}_{oppo}$ & $0$ ft/s & $2200$ ft/s \\
%         & $\theta_{bea, oppo}, \theta_{ele, oppo}$ & $-\pi$ rad & $\pi$ rad \\
%         \midrule
%         \textbf{Historical Actions} 
%         & $\delta_a, \delta_e, \delta_r, \delta_T$ & $-1$ & $1$ \\
%         \botrule
%     \end{tabular}
% \end{table*}




%%=============================================%%
%% For submissions to Nature Portfolio Journals %%
%% please use the heading ``Extended Data''.   %%
%%=============================================%%

%%=============================================================%%
%% Sample for another appendix section			       %%
%%=============================================================%%

%% \section{Example of another appendix section}\label{secA2}%
%% Appendices may be used for helpful, supporting or essential material that would otherwise 
%% clutter, break up or be distracting to the text. Appendices can consist of sections, figures, 
%% tables and equations etc.

\end{appendices}

%%===========================================================================================%%
%% If you are submitting to one of the Nature Portfolio journals, using the eJP submission   %%
%% system, please include the references within the manuscript file itself. You may do this  %%
%% by copying the reference list from your .bbl file, paste it into the main manuscript .tex %%
%% file, and delete the associated \verb+\bibliography+ commands.                            %%
%%===========================================================================================%%

\bibliography{sn-bibliography}% common bib file
%% if required, the content of .bbl file can be included here once bbl is generated
%%\input sn-article.bbl

\end{document}
